%% toc: P2 23 23.1,23.2,23.3,23.4
\section{Proofs for Church's type theory, full categoricity and nonaxiomatizability}
\label{sec:annex-p02-w02-p02-040}

Notation for Church's system is that of
\cref{def:p02-w02-p02-040-church-types,def:p02-w02-p02-040-church-wff,def:p02-w02-p02-040-church-abbrev,def:p02-w02-p02-040-church-rules,def:p02-w02-p02-040-church-axioms}.
The propositional input (P) of \cref{def:p02-w02-p02-040-church-prop},
proved in \cref{annex:p02-w02-p02-040-church-prop-complete} below without
the deduction theorem, is used as follows: if \(\tau(p_1,\ldots,p_k)\) is a tautology in \(\sim,\vee\) and
\(\mathbf B^1_o,\ldots,\mathbf B^k_o\) are wffs, then
\(\derives\tau(\mathbf B^1_o,\ldots,\mathbf B^k_o)\).  We write
\emph{(P)} for an appeal to it, and \emph{(P), V} when the tautology is an
implication chain \(\mathbf C^1\supset.\,\cdots\supset.\,\mathbf C^m\supset\mathbf D\)
whose antecedents are already available, so that \(m\) applications of
Rule~V give \(\mathbf D\).

\paragraph{Renaming convention.}  Rule~I renames a bound variable to one that
does not occur in the part concerned.  Before an application of IV\('\),
II or III whose proviso about bound variables would fail, we rename the
offending bound variables to variables occurring nowhere in the proof, and
after it we rename them back.  Renaming back is possible: each restored
binder heads a part that the intervening step copied unchanged, and the
original name does not occur in that part, since inside a part
\((\lambda\mathbf y\mathbf M)\) every occurrence of \(\mathbf y\) is bound.
Abbreviations such as \(Q_{o\alpha\alpha}\), \(N_{o\alpha'}\) and
\(S_{\alpha'\alpha'}\) are always taken in a variant whose bound variables
are distinct from all other variables in play.  We do not mention these
renamings again.

\begin{annexitem}{annex:p02-w02-p02-040-church-prop-complete}{thm:p02-w02-p02-040-church-prop-complete}{Axioms 1--4 yield every tautology instance}
\begin{proof}
Bold letters \(\mathbf A,\mathbf B,\mathbf C\) denote wffs of type \(o\), and
\(\mathbf A\supset\mathbf B\) is the formula \(\sim\mathbf A\vee\mathbf B\).
The route below reduces a tautologous disjunction to shorter ones; it is
not the route of the works the source cites.

\emph{Step 1: axiom instances.}  For all \(\mathbf A,\mathbf B,\mathbf C\) the
following are theorems:
\[
 \text{(a1)}\ \mathbf A\vee\mathbf A\supset\mathbf A,\quad
 \text{(a2)}\ \mathbf A\supset\mathbf A\vee\mathbf B,\quad
 \text{(a3)}\ \mathbf A\vee\mathbf B\supset\mathbf B\vee\mathbf A,\quad
 \text{(a4)}\ \mathbf A\supset\mathbf B\supset.\,\mathbf C\vee\mathbf A\supset\mathbf C\vee\mathbf B.
\]
For (a4), choose variables \(p^*,q^*,r^*\) of type \(o\) occurring in none of
\(\mathbf A,\mathbf B,\mathbf C\).  Axiom~4 has no bound variable, so IV\('\)
substitutes \(p^*,q^*,r^*\) for \(p,q,r\) one at a time.  Then IV\('\)
substitutes \(\mathbf A\) for \(p^*\), \(\mathbf B\) for \(q^*\) and
\(\mathbf C\) for \(r^*\), with the renaming convention whenever a bound
variable of a wff already inserted is free in the next one.  No substituted
variable occurs in a wff inserted earlier, so the result is (a4).  The same
procedure applied to Axioms~1--3 gives (a1)--(a3).

\emph{Step 2: schematic derivations.}  Let \(P_1,P_2,\ldots\) be letters, and
call the formulas built from letters by \(\sim\) and \(\vee\)
\emph{propositional formulas} (\(\supset\) again abbreviating).  A
\emph{schematic derivation} is a finite sequence of propositional formulas,
each an instance of (a1)--(a4) at propositional formulas or obtained from
two earlier members \(X\) and \(X\supset Y\) as \(Y\).  Write \(\derives_0X\)
if \(X\) ends one.  If \(\sigma\) assigns a wff of type \(o\) to each letter
and \(X_1,\ldots,X_m\) is a schematic derivation, then each \(\sigma X_i\) is
an instance of (a1)--(a4) at wffs, hence a theorem by Step~1, or follows from
earlier members by Rule~V; so \(\derives\sigma X_m\).  It therefore suffices to
prove \(\derives_0X\) for every tautologous propositional formula \(X\): for
\(\tau(p_1,\ldots,p_k)\), apply this to the formula \(\hat\tau\) with \(P_i\)
in place of \(p_i\) and take \(\sigma(P_i)=\mathbf B^i_o\).  From here on
\(X,Y,Z,R,S\) are propositional formulas.

\emph{Step 3: derived rules and theorems.}  Each line is obtained by Rule~V
from the instances named.
\begin{enumerate}
 \item[(r1)] From \(\derives_0X\supset Y\) and \(\derives_0Y\supset Z\),
   \(\derives_0X\supset Z\): (a4) with \(Y,Z,\sim X\) for \(\mathbf A,\mathbf B,\mathbf C\)
   is \((Y\supset Z)\supset((X\supset Y)\supset(X\supset Z))\).
 \item[(r2)] From \(\derives_0X\supset Y\), \(\derives_0Z\vee X\supset Z\vee Y\), by (a4).
 \item[(r3)] From \(\derives_0X\supset Y\), \(\derives_0X\vee Z\supset Y\vee Z\):
   (a3), then (r2), then (a3), joined by (r1).
 \item[(r4)] From \(\derives_0X\supset Z\) and \(\derives_0Y\supset Z\),
   \(\derives_0X\vee Y\supset Z\): (r2) gives \(X\vee Y\supset X\vee Z\), (r3)
   gives \(X\vee Z\supset Z\vee Z\), and (a1) gives \(Z\vee Z\supset Z\).
 \item[(t1)] \(\derives_0X\supset X\), that is \(\derives_0\sim X\vee X\): (a2)
   gives \(X\supset X\vee X\), and (a1) with (r1) finishes.
 \item[(t2)] \(\derives_0Y\supset X\vee Y\): (a2) gives \(Y\supset Y\vee X\), and
   (a3) with (r1) finishes.
 \item[(t3)] \(\derives_0X\supset\sim\sim X\): (t1) at \(\sim X\) is
   \(\sim\sim X\vee\sim X\), and (a3) turns it into \(\sim X\vee\sim\sim X\).
 \item[(r5)] (Cut.)  From \(\derives_0R\vee X\) and \(\derives_0\sim X\vee S\),
   \(\derives_0R\vee S\): the second premise is \(X\supset S\); (r2) gives
   \(R\vee X\supset R\vee S\).
\end{enumerate}

\emph{Step 4: disjunctions.}  A \emph{reading} of a formula as a disjunction
tree is defined by: any formula \(X\) may be read as a tree with the single
leaf \(X\); if \(X_1\) and \(X_2\) are read as trees, \(X_1\vee X_2\) may be
read as the tree whose leaf list is that of \(X_1\) followed by that of
\(X_2\).  \emph{Subsumption:} if \(X\) and \(Y\) are read as trees and every
leaf of \(X\) is a leaf of \(Y\), then \(\derives_0X\supset Y\).  First, if
\(L\) is a leaf of \(Y\), then \(\derives_0L\supset Y\), by induction on the
reading of \(Y\): for a single leaf use (t1); for \(Y_1\vee Y_2\) with \(L\)
in \(Y_1\) use the hypothesis, (a2) and (r1), and with \(L\) in \(Y_2\) the
hypothesis, (t2) and (r1).  Then induct on the reading of \(X\), using (r4)
at \(X_1\vee X_2\).

For a nonempty list \(\Gamma=(X_1,\ldots,X_n)\) put \(\bigvee\Gamma=X_1\) if
\(n=1\) and \(\bigvee\Gamma=X_1\vee\bigvee(X_2,\ldots,X_n)\) otherwise, read
with leaves \(X_1,\ldots,X_n\).  \emph{Replacement:} if \(\derives_0Y\supset Y'\)
and \(\Gamma'\) is \(\Gamma\) with one entry \(Y\) replaced by \(Y'\), then
\(\derives_0\bigvee\Gamma\supset\bigvee\Gamma'\).  For \(n=1\) this is the
hypothesis.  Otherwise let \(R\) be \(\bigvee\) of the other entries in order;
subsumption gives \(\bigvee\Gamma\supset Y\vee R\), (r3) gives
\(Y\vee R\supset Y'\vee R\), subsumption gives \(Y'\vee R\supset\bigvee\Gamma'\),
and (r1) joins them.

\emph{Step 5: tautologous disjunctions are derivable.}  A \emph{literal} is a
letter or the negation of a letter.  Let \(w(X)\) be the number of
occurrences of \(\sim\) and \(\vee\) in \(X\), and for a list \(\Gamma\) let
\(\mu(\Gamma)\) be the sum of \(w(X_i)\) over the entries that are not
literals.  We show by induction on \(\mu(\Gamma)\) that if \(\bigvee\Gamma\) is
a tautology then \(\derives_0\bigvee\Gamma\).

If \(\mu(\Gamma)=0\), every entry is a literal.  Some letter \(P\) has both
\(P\) and \(\sim P\) among the entries: otherwise the valuation making each
letter entry false and each negated-letter entry true (and other letters
arbitrary) makes \(\bigvee\Gamma\) false.  By (t1), \(\derives_0\sim P\vee P\),
and subsumption gives \(\derives_0\bigvee\Gamma\).

If \(\mu(\Gamma)>0\), take the first entry \(X_i\) that is not a literal.  It
has one of three forms.
\begin{enumerate}
 \item \(X_i=Y\vee Z\).  Let \(\Gamma'\) replace this entry by the two entries
   \(Y,Z\).  Then \(\bigvee\Gamma'\) has the truth value of \(\bigvee\Gamma\)
   under every valuation and \(\mu(\Gamma')\le\mu(\Gamma)-1\), so
   \(\derives_0\bigvee\Gamma'\).  Read \(\bigvee\Gamma\) with the leaf
   \(Y\vee Z\) split into leaves \(Y\) and \(Z\); subsumption gives
   \(\bigvee\Gamma'\supset\bigvee\Gamma\), and Rule~V gives \(\bigvee\Gamma\).
 \item \(X_i=\sim\sim Y\).  Let \(\Gamma'\) replace it by \(Y\).  Truth values
   agree and \(\mu\) drops by at least two, so \(\derives_0\bigvee\Gamma'\);
   (t3) and replacement give \(\bigvee\Gamma'\supset\bigvee\Gamma\).
 \item \(X_i=\sim(Y\vee Z)\).  Let \(\Gamma_1\) and \(\Gamma_2\) replace it by
   \(\sim Y\) and by \(\sim Z\).  Under a valuation making \(\bigvee\Gamma\)
   true, either another entry is true or \(\sim(Y\vee Z)\) is true, and then
   \(\sim Y\) and \(\sim Z\) are both true; so both are tautologies, and since
   \(w(\sim Y)<w(\sim(Y\vee Z))\), both are derivable.  Next, (t1) at
   \(Y\vee Z\) gives \(\sim(Y\vee Z)\vee(Y\vee Z)\); subsumption gives
   \(\bigvee(\sim(Y\vee Z),Y,Z)\); replacement twice with (t3) gives
   \(\bigvee(\sim(Y\vee Z),\sim\sim Y,\sim\sim Z)\); and subsumption gives
   \[
    E:\quad\sim\sim Y\vee(\sim\sim Z\vee\sim(Y\vee Z)).
   \]
   If \(n=1\), \(E\) is \(\sim Y\supset.\,\sim Z\supset\sim(Y\vee Z)\), and
   Rule~V twice with \(\derives_0\sim Y\) and \(\derives_0\sim Z\) gives
   \(\sim(Y\vee Z)\).  If \(n\ge2\), let \(R\) be \(\bigvee\) of the other entries.
   Subsumption turns \(\bigvee\Gamma_1\) into \(R\vee\sim Y\); the cut (r5)
   with \(E\) gives \(R\vee(\sim\sim Z\vee\sim(Y\vee Z))\), which subsumption
   turns into \(\sim\sim Z\vee(R\vee\sim(Y\vee Z))\).  Subsumption turns
   \(\bigvee\Gamma_2\) into \(R\vee\sim Z\), and the cut gives
   \(R\vee(R\vee\sim(Y\vee Z))\).  Every leaf of this formula, read with the
   leaves of \(R\), is a leaf of \(\bigvee\Gamma\), so subsumption gives
   \(\derives_0\bigvee\Gamma\).
\end{enumerate}

\emph{Step 6.}  A tautology \(X\) is \(\bigvee(X)\), so \(\derives_0X\) by
Step~5, and Step~2 transfers this to every substitution instance.  The rules
used are I and IV\('\) (Step~1) and V.
\end{proof}
\end{annexitem}

\begin{annexitem}{annex:p02-w02-p02-040-church-deduction}{thm:p02-w02-p02-040-church-deduction}{Deduction theorem (Church's VII)}
\begin{proof}
\emph{Step 1: theorems may be used under assumptions.}  Suppose
\(\derives\mathbf B_o\) and let \(\mathbf A^1_o,\ldots,\mathbf A^n_o\) be
given.  Let \(\mathbf y^1,\ldots,\mathbf y^k\) be the variables that occur
in a fixed proof of \(\mathbf B_o\) and are free in some
\(\mathbf A^i_o\), and pick distinct variables \(\mathbf z^1,\ldots,\mathbf z^k\)
of the same types occurring neither in that proof nor in any assumption.
Replacing each \(\mathbf y^j\) by \(\mathbf z^j\) everywhere in the proof
gives a proof: the replacement is injective and type-preserving, so it
carries variants of axioms to variants of axioms, and it preserves every
proviso of Rules I--VI, all of which speak only of occurrence and
distinctness of variables.  The new proof mentions no \(\mathbf y^j\), so it
is a proof on the assumptions.  Its last line \(\mathbf B'_o\) is
\(\mathbf B_o\) with each \(\mathbf y^j\) replaced by \(\mathbf z^j\).
Apply IV\('\) to substitute \(\mathbf y^j\) for the free occurrences of
\(\mathbf z^j\), one \(j\) at a time; this substitutes for the
\(\mathbf z^j\), which are not free in any assumption, and its proviso holds
because \(\mathbf y^j\) does not occur in \(\mathbf B'_o\).  Then Rule~I
renames each bound \(\mathbf z^j\) back to \(\mathbf y^j\).  The result is
\(\mathbf B_o\), so \(\mathbf A^1_o,\ldots,\mathbf A^n_o\derives\mathbf B_o\).

\emph{Step 2: the deduction theorem.}  Let
\(\mathbf B^1_o,\ldots,\mathbf B^m_o\) be a proof of \(\mathbf B_o\) on
\(\mathbf A^1_o,\ldots,\mathbf A^n_o\), with \(\mathbf B^m_o=\mathbf B_o\).
Write \(\Gamma\) for \(\mathbf A^1_o,\ldots,\mathbf A^{n-1}_o\).  We show by
induction on \(i\) that
\(\Gamma\derives\mathbf A^n_o\supset\mathbf B^i_o\) for \(i=1,\ldots,m\),
following the cases of \cite[\S5, p.~63]{Church1940}.
\begin{enumerate}
 \item \(\mathbf B^i_o\) is \(\mathbf A^n_o\).  The formula
   \(\mathbf A^n_o\supset\mathbf A^n_o\) is a theorem by (P), hence derivable
   from \(\Gamma\) by Step~1.
 \item \(\mathbf B^i_o\) is some \(\mathbf A^j_o\) with \(j<n\), or a variant of
   an axiom.  Then \(\Gamma\derives\mathbf B^i_o\) in one line, and
   \(\mathbf B^i_o\supset.\,\mathbf A^n_o\supset\mathbf B^i_o\) is a theorem by
   (P); Step~1 and Rule~V give the claim.
 \item \(\mathbf B^i_o\) comes from \(\mathbf B^a_o\), \(a<i\), by one of
   Rules I--III.  The part replaced in \(\mathbf B^a_o\) is also a part of
   \(\mathbf A^n_o\supset\mathbf B^a_o\), which abbreviates
   \(A_{ooo}(N_{oo}\mathbf A^n_o)\mathbf B^a_o\); the same replacement
   there yields \(\mathbf A^n_o\supset\mathbf B^i_o\).  Conversion rules carry
   no restriction in proofs on assumptions.
 \item \(\mathbf B^i_o\) is \(\mathbf F_{o\alpha}\mathbf A_\alpha\) and comes by
   Rule~IV from \(\mathbf B^a_o=\mathbf F_{o\alpha}\mathbf x_\alpha\), with
   \(\mathbf x_\alpha\) not free in \(\mathbf F_{o\alpha}\).  The restriction on
   proofs on assumptions makes \(\mathbf x_\alpha\) not free in any
   \(\mathbf A^j_o\), in particular not in \(\mathbf A^n_o\).  Apply IV\('\)
   to \(\Gamma\derives\mathbf A^n_o\supset\mathbf F_{o\alpha}\mathbf x_\alpha\),
   substituting \(\mathbf A_\alpha\) for \(\mathbf x_\alpha\).  The only free
   occurrence of \(\mathbf x_\alpha\) is the displayed argument, so the result
   is \(\mathbf A^n_o\supset\mathbf F_{o\alpha}\mathbf A_\alpha\), and the
   substituted variable is not free in \(\Gamma\).
 \item \(\mathbf B^i_o\) comes by Rule~V from \(\mathbf B^a_o\) and
   \(\mathbf B^b_o=\mathbf B^a_o\supset\mathbf B^i_o\), \(a,b<i\).  From the
   induction hypotheses \(\mathbf A^n_o\supset\mathbf B^a_o\) and
   \(\mathbf A^n_o\supset.\,\mathbf B^a_o\supset\mathbf B^i_o\), (P), V with
   the tautology
   \((p\supset q)\supset.\,(p\supset.\,q\supset r)\supset.\,p\supset r\) give
   \(\mathbf A^n_o\supset\mathbf B^i_o\).
 \item \(\mathbf B^i_o\) is \(\Pi_{o(o\alpha)}\mathbf F_{o\alpha}\) and comes by
   Rule~VI from \(\mathbf B^a_o=\mathbf F_{o\alpha}\mathbf x_\alpha\), with
   \(\mathbf x_\alpha\) not free in \(\mathbf F_{o\alpha}\) and, by the
   restriction, not free in any assumption.  By the induction hypothesis and
   VI\('\), which generalizes upon \(\mathbf x_\alpha\),
   \(\Gamma\derives(\mathbf x_\alpha)[\sim\mathbf A^n_o\vee\mathbf F_{o\alpha}\mathbf x_\alpha]\).
   Substituting \(\sim\mathbf A^n_o\) for \(p\) and \(\mathbf F_{o\alpha}\) for
   \(f_{o\alpha}\) in Axiom 6\(^\alpha\) by IV\('\) gives the theorem
   \((\mathbf x_\alpha)[\sim\mathbf A^n_o\vee\mathbf F_{o\alpha}\mathbf x_\alpha]
   \supset\sim\mathbf A^n_o\vee\Pi_{o(o\alpha)}\mathbf F_{o\alpha}\); the
   proviso holds since \(\mathbf x_\alpha\) is free in neither
   \(\mathbf A^n_o\) nor \(\mathbf F_{o\alpha}\).  Step~1 and Rule~V give
   \(\Gamma\derives\mathbf A^n_o\supset\Pi_{o(o\alpha)}\mathbf F_{o\alpha}\).
\end{enumerate}
The new derivations substitute for, or generalize upon, only variables
that the original proof already substituted for or generalized upon, and
variables of inserted theorems, which Step~1 handles; so each is a proof
on \(\Gamma\).  The case \(i=m\) is the theorem.
\end{proof}
\end{annexitem}

\begin{annexitem}{annex:p02-w02-p02-040-church-12-21}{thm:p02-w02-p02-040-church-12-21}{Quantifier and equality theorems 12--21}
\begin{proof}
Recall that \((\mathbf x_\alpha)\mathbf A_o\) is
\(\Pi_{o(o\alpha)}(\lambda\mathbf x_\alpha\mathbf A_o)\) and
\((\exists\mathbf x_\alpha)\mathbf A_o\) is
\(\sim\Pi_{o(o\alpha)}(\lambda\mathbf x_\alpha[\sim\mathbf A_o])\).

\emph{12\(^\alpha\).}  In 5\(^\alpha\) substitute \(y_\alpha\) for
\(x_\alpha\) and then \(\lambda x_\alpha(f_{o\alpha}x_\alpha)\) for
\(f_{o\alpha}\) (IV\('\) twice), obtaining
\(\Pi_{o(o\alpha)}(\lambda x_\alpha(f_{o\alpha}x_\alpha))\supset
(\lambda x_\alpha(f_{o\alpha}x_\alpha))y_\alpha\).  Rule~II on the last part
gives 12\(^\alpha\).

\emph{13\(^\alpha\).}  In 12\(^\alpha\) substitute
\(\lambda z_\alpha[\sim f_{o\alpha}z_\alpha]\) for \(f_{o\alpha}\) and reduce
both new redexes by Rule~II:
\(\Pi_{o(o\alpha)}(\lambda x_\alpha[\sim f_{o\alpha}x_\alpha])\supset\sim f_{o\alpha}y_\alpha\).
(P), V with \((p\supset\sim q)\supset.\,q\supset\sim p\) gives
\(f_{o\alpha}y_\alpha\supset\sim\Pi_{o(o\alpha)}(\lambda x_\alpha[\sim f_{o\alpha}x_\alpha])\),
which is 13\(^\alpha\).

\emph{14\(^\alpha\).}  In 6\(^\alpha\) substitute \(\sim p\) for \(p\); the
proviso holds since \(x_\alpha\) is not free in \(\sim p\).  Because
\([\mathbf A\supset\mathbf B]\) abbreviates \([[\sim\mathbf A]\vee\mathbf B]\),
the result is
\((x_\alpha)[p\supset f_{o\alpha}x_\alpha]\supset.\,p\supset\Pi_{o(o\alpha)}f_{o\alpha}\).
Substitute \(\lambda z_\alpha(f_{o\alpha}z_\alpha)\) for \(f_{o\alpha}\), reduce
the redex \((\lambda z_\alpha(f_{o\alpha}z_\alpha))x_\alpha\) by Rule~II, and
rename \(z_\alpha\) to \(x_\alpha\) in the consequent by Rule~I.  This is
14\(^\alpha\).

\emph{Derived rule G.}  If \(\derives\mathbf A_o\supset\mathbf B_o\) and
\(\mathbf x_\alpha\) is not free in \(\mathbf A_o\), then
\(\derives\mathbf A_o\supset(\mathbf x_\alpha)\mathbf B_o\).  Indeed VI\('\)
gives \((\mathbf x_\alpha)[\mathbf A_o\supset\mathbf B_o]\); substituting
\(\mathbf A_o\) for \(p\) and \(\lambda\mathbf x_\alpha\mathbf B_o\) for
\(f_{o\alpha}\) in 14\(^\alpha\) and reducing gives
\((\mathbf x_\alpha)[\mathbf A_o\supset\mathbf B_o]\supset.\,\mathbf A_o\supset(\mathbf x_\alpha)\mathbf B_o\),
the proviso holding because \(\mathbf x_\alpha\) is not free in
\(\mathbf A_o\); Rule~V finishes.  The same argument under assumptions
works when \(\mathbf x_\alpha\) is also not free in any assumption.

\emph{15\(^\alpha\).}  Substituting \(\lambda x_\alpha[f_{o\alpha}x_\alpha\supset p]\)
for \(f_{o\alpha}\) and \(x_\alpha\) for \(y_\alpha\) in 12\(^\alpha\) and
reducing gives
\((x_\alpha)[f_{o\alpha}x_\alpha\supset p]\supset.\,f_{o\alpha}x_\alpha\supset p\).
(P), V with \((q\supset.\,r\supset p)\supset.\,q\supset.\,\sim p\supset\sim r\)
give \((x_\alpha)[f_{o\alpha}x_\alpha\supset p]\supset.\,\sim p\supset\sim f_{o\alpha}x_\alpha\),
and G (the antecedent has no free \(x_\alpha\)) gives
\((x_\alpha)[f_{o\alpha}x_\alpha\supset p]\supset(x_\alpha)[\sim p\supset\sim f_{o\alpha}x_\alpha]\).
Substituting \(\sim p\) for \(p\) and \(\lambda x_\alpha[\sim f_{o\alpha}x_\alpha]\)
for \(f_{o\alpha}\) in 14\(^\alpha\) and reducing gives
\((x_\alpha)[\sim p\supset\sim f_{o\alpha}x_\alpha]\supset.\,\sim p\supset(x_\alpha)[\sim f_{o\alpha}x_\alpha]\).
(P), V with \((a\supset b)\supset.\,(b\supset.\,\sim p\supset c)\supset.\,a\supset.\,\sim c\supset p\)
then give
\((x_\alpha)[f_{o\alpha}x_\alpha\supset p]\supset.\,\sim(x_\alpha)[\sim f_{o\alpha}x_\alpha]\supset p\),
which is 15\(^\alpha\).

\emph{Derived rule E.}  If \(\derives\mathbf B_o\supset\mathbf A_o\) and
\(\mathbf x_\alpha\) is not free in \(\mathbf A_o\), then
\(\derives(\exists\mathbf x_\alpha)\mathbf B_o\supset\mathbf A_o\): apply VI\('\),
then 15\(^\alpha\) with \(\lambda\mathbf x_\alpha\mathbf B_o\) for
\(f_{o\alpha}\) and \(\mathbf A_o\) for \(p\), and Rule~V.

\emph{16\(^\alpha\).}  By (P), \(\derives f_{o\alpha}x_\alpha\supset f_{o\alpha}x_\alpha\).
VI\('\) gives \((f_{o\alpha})[f_{o\alpha}x_\alpha\supset f_{o\alpha}x_\alpha]\).
Rule~III twice (the bound variable \(f_{o\alpha}\) differs from
\(x_\alpha\)) rewrites this as
\((\lambda u_\alpha\lambda v_\alpha[(f_{o\alpha})[f_{o\alpha}u_\alpha\supset f_{o\alpha}v_\alpha]])x_\alpha x_\alpha\),
which is \(Q_{o\alpha\alpha}x_\alpha x_\alpha\), that is \(x_\alpha=x_\alpha\).
Only Axioms 1--4 and the rules were used.

\emph{17\(^\alpha\).}  By Rule~II twice, \(x_\alpha=y_\alpha\) converts to
\((g_{o\alpha})[g_{o\alpha}x_\alpha\supset g_{o\alpha}y_\alpha]\).  Apply
12 at type \(o\alpha\), substituting
\(\lambda g_{o\alpha}[g_{o\alpha}x_\alpha\supset g_{o\alpha}y_\alpha]\) for its
function variable and \(f_{o\alpha}\) for its instance variable, and reduce:
\((g_{o\alpha})[g_{o\alpha}x_\alpha\supset g_{o\alpha}y_\alpha]\supset.\,f_{o\alpha}x_\alpha\supset f_{o\alpha}y_\alpha\).
Rule~III on the antecedent restores \(x_\alpha=y_\alpha\).

\emph{18\(^{\beta\alpha}\).}  Substitute
\(\lambda z_\alpha[f_{\beta\alpha}x_\alpha=f_{\beta\alpha}z_\alpha]\) for
\(f_{o\alpha}\) in 17\(^\alpha\) and reduce:
\(x_\alpha=y_\alpha\supset.\,f_{\beta\alpha}x_\alpha=f_{\beta\alpha}x_\alpha\supset f_{\beta\alpha}x_\alpha=f_{\beta\alpha}y_\alpha\).
16\(^\beta\) with \(f_{\beta\alpha}x_\alpha\) for \(x_\beta\) gives the middle
formula, and (P), V with
\((p\supset.\,q\supset r)\supset.\,q\supset.\,p\supset r\) give 18\(^{\beta\alpha}\).

\emph{19\(^\alpha\).}  Substitute \(\lambda z_\alpha[z_\alpha=x_\alpha]\) for
\(f_{o\alpha}\) in 17\(^\alpha\) and reduce:
\(x_\alpha=y_\alpha\supset.\,x_\alpha=x_\alpha\supset y_\alpha=x_\alpha\).
Conclude with 16\(^\alpha\) and the tautology used for 18.

\emph{20\(^\alpha\).}  In 17\(^\alpha\) substitute \(z_\alpha\) for \(y_\alpha\),
then \(y_\alpha\) for \(x_\alpha\), then
\(\lambda w_\alpha[x_\alpha=w_\alpha]\) for \(f_{o\alpha}\), and reduce:
\(y_\alpha=z_\alpha\supset.\,x_\alpha=y_\alpha\supset x_\alpha=z_\alpha\).
(P), V with \((p\supset.\,q\supset r)\supset.\,q\supset.\,p\supset r\)
give 20\(^\alpha\).

\emph{21\(^{\alpha\beta}\).}  In 10\(^{\alpha\beta}\) substitute
\(\lambda z_\beta(f_{\alpha\beta}z_\beta)\) for \(g_{\alpha\beta}\) and reduce:
\((x_\beta)[f_{\alpha\beta}x_\beta=f_{\alpha\beta}x_\beta]\supset f_{\alpha\beta}=\lambda z_\beta(f_{\alpha\beta}z_\beta)\).
16\(^\alpha\) with \(f_{\alpha\beta}x_\beta\) for \(x_\alpha\) and VI\('\) give the
antecedent; Rule~V and a renaming of \(z_\beta\) to \(x_\beta\) give
21\(^{\alpha\beta}\).  Only Axioms 1--4 (through 16) and 10\(^{\alpha\beta}\)
were used.
\end{proof}
\end{annexitem}

\begin{annexitem}{annex:p02-w02-p02-040-church-comprehension}{thm:p02-w02-p02-040-church-comprehension}{Comprehension by abstraction}
\begin{proof}
By (P), \(\derives\mathbf A_o\equiv\mathbf A_o\).  Rule~III replaces the first
occurrence of \(\mathbf A_o\) by \((\lambda\mathbf x_\alpha\mathbf A_o)\mathbf x_\alpha\),
since substituting \(\mathbf x_\alpha\) for \(\mathbf x_\alpha\) in
\(\mathbf A_o\) returns \(\mathbf A_o\).  VI\('\) gives
\((\mathbf x_\alpha)[(\lambda\mathbf x_\alpha\mathbf A_o)\mathbf x_\alpha\equiv\mathbf A_o]\).
Now apply 13 at type \(o\alpha\): substitute
\(\lambda\mathbf f_{o\alpha}(\mathbf x_\alpha)[\mathbf f_{o\alpha}\mathbf x_\alpha\equiv\mathbf A_o]\)
for its function variable (of type \(o(o\alpha)\)) and
\(\lambda\mathbf x_\alpha\mathbf A_o\) for its instance variable, and reduce.
Since \(\mathbf f_{o\alpha}\) is not free in \(\mathbf A_o\), the
reduced antecedent is the formula just proved and the reduced consequent is
\((\exists\mathbf f_{o\alpha})(\mathbf x_\alpha)[\mathbf f_{o\alpha}\mathbf x_\alpha\equiv\mathbf A_o]\).
Rule~V gives the theorem.  The axioms used are 1--4 and, through 13, 5.
\end{proof}
\end{annexitem}

\begin{annexitem}{annex:p02-w02-p02-040-church-induction}{thm:p02-w02-p02-040-church-induction}{Peano theorems 22--24 and the induction theorem}
\begin{proof}
Write \(H(\mathbf g)\) for
\((w_{\alpha'})[\mathbf g\,w_{\alpha'}\supset\mathbf g(S_{\alpha'\alpha'}w_{\alpha'})]\)
with a bound variable \(w_{\alpha'}\) chosen fresh.  By Rule~II,
\(N_{o\alpha'}\mathbf A_{\alpha'}\) converts to
\((f_{o\alpha'})[f_{o\alpha'}0_{\alpha'}\supset.\,H(f_{o\alpha'})\supset f_{o\alpha'}\mathbf A_{\alpha'}]\),
and conversely by Rule~III.

\emph{22\(^\alpha\).}  (P) with \(p\supset.\,q\supset p\) gives
\(f_{o\alpha'}0_{\alpha'}\supset.\,H(f_{o\alpha'})\supset f_{o\alpha'}0_{\alpha'}\);
VI\('\) on \(f_{o\alpha'}\) and Rule~III give \(N_{o\alpha'}0_{\alpha'}\).

\emph{23\(^\alpha\).}  Assume \(N_{o\alpha'}x_{\alpha'}\) and convert it.
Twelve at type \(o\alpha'\), instantiated at \(f_{o\alpha'}\) itself, and
Rule~V give
\(f_{o\alpha'}0_{\alpha'}\supset.\,H(f_{o\alpha'})\supset f_{o\alpha'}x_{\alpha'}\).
Twelve at type \(\alpha'\) with
\(\lambda w_{\alpha'}[f_{o\alpha'}w_{\alpha'}\supset f_{o\alpha'}(S_{\alpha'\alpha'}w_{\alpha'})]\)
for the function variable and \(x_{\alpha'}\) for the instance gives the
theorem
\(H(f_{o\alpha'})\supset.\,f_{o\alpha'}x_{\alpha'}\supset f_{o\alpha'}(S_{\alpha'\alpha'}x_{\alpha'})\).
(P), V with
\((a\supset.\,b\supset c)\supset.\,(b\supset.\,c\supset d)\supset.\,a\supset.\,b\supset d\)
give
\(f_{o\alpha'}0_{\alpha'}\supset.\,H(f_{o\alpha'})\supset f_{o\alpha'}(S_{\alpha'\alpha'}x_{\alpha'})\)
on the assumption.  The variable \(f_{o\alpha'}\) is bound in the
assumption, so VI\('\) may generalize upon it; Rule~III then gives
\(N_{o\alpha'}(S_{\alpha'\alpha'}x_{\alpha'})\).  The deduction theorem gives
23\(^\alpha\).

\emph{24\(^\alpha\).}  Let \(\mathbf g\) be
\(\lambda x_{\alpha'}[N_{o\alpha'}x_{\alpha'}\,f_{o\alpha'}x_{\alpha'}]\), the
juxtaposition inside the brackets being conjunction.  Take the assumptions
\(\mathbf A^1=f_{o\alpha'}0_{\alpha'}\),
\(\mathbf A^2=(x_{\alpha'})[N_{o\alpha'}x_{\alpha'}\supset.\,f_{o\alpha'}x_{\alpha'}\supset f_{o\alpha'}(S_{\alpha'\alpha'}x_{\alpha'})]\)
and \(\mathbf A^3=N_{o\alpha'}x_{\alpha'}\).  Choose \(w_{\alpha'}\) different
from \(x_{\alpha'}\) and not free in any assumption.
\begin{enumerate}
 \item From 22\(^\alpha\), \(\mathbf A^1\), (P), V with
   \(p\supset.\,q\supset pq\), and Rule~III:
   \(\mathbf A^1\derives\mathbf g\,0_{\alpha'}\).
 \item Twelve with \(w_{\alpha'}\) for the instance and Rule~V give
   \(\mathbf A^2\derives N_{o\alpha'}w_{\alpha'}\supset.\,f_{o\alpha'}w_{\alpha'}\supset f_{o\alpha'}(S_{\alpha'\alpha'}w_{\alpha'})\).
   With 23\(^\alpha\) at \(w_{\alpha'}\), (P), V with
   \((a\supset.\,b\supset c)\supset.\,(a\supset d)\supset.\,ab\supset dc\) give
   \[
    \mathbf A^2\derives[N_{o\alpha'}w_{\alpha'}\,f_{o\alpha'}w_{\alpha'}]\supset
    [N_{o\alpha'}(S_{\alpha'\alpha'}w_{\alpha'})\,f_{o\alpha'}(S_{\alpha'\alpha'}w_{\alpha'})];
   \]
   Rule~III gives \(\mathbf g\,w_{\alpha'}\supset\mathbf g(S_{\alpha'\alpha'}w_{\alpha'})\),
   and VI\('\) on \(w_{\alpha'}\) gives \(\mathbf A^2\derives H(\mathbf g)\).
 \item Converting \(\mathbf A^3\) and applying 12 at type \(o\alpha'\) with
   \(\mathbf g\) for the instance give
   \(\mathbf A^3\derives\mathbf g\,0_{\alpha'}\supset.\,H(\mathbf g)\supset\mathbf g\,x_{\alpha'}\).
\end{enumerate}
Rule~V twice gives \(\mathbf A^1,\mathbf A^2,\mathbf A^3\derives\mathbf g\,x_{\alpha'}\);
Rule~II and (P), V with \(pq\supset q\) give
\(f_{o\alpha'}x_{\alpha'}\).  Only \(w_{\alpha'}\) was generalized upon and
no assumption variable was substituted for, so the deduction theorem applies
three times and yields 24\(^\alpha\).

\emph{VIII.}  From the two hypotheses, the deduction theorem gives
\(\mathbf A^1_o,\ldots,\mathbf A^n_o\derives N_{o\alpha'}\mathbf x_{\alpha'}\supset.\,
\mathbf F_{o\alpha'}\mathbf x_{\alpha'}\supset\mathbf F_{o\alpha'}(S_{\alpha'\alpha'}\mathbf x_{\alpha'})\),
and VI\('\), legitimate because \(\mathbf x_{\alpha'}\) is not free in the
assumptions, gives its generalization on \(\mathbf x_{\alpha'}\).
Substituting \(\mathbf F_{o\alpha'}\) for \(f_{o\alpha'}\) in 24\(^\alpha\),
written with \(\mathbf x_{\alpha'}\) as its variable, is permitted because
\(\mathbf x_{\alpha'}\) is not free in \(\mathbf F_{o\alpha'}\); it gives
\(\mathbf F_{o\alpha'}0_{\alpha'}\supset.\,(\mathbf x_{\alpha'})[N_{o\alpha'}\mathbf x_{\alpha'}\supset.\,
\mathbf F_{o\alpha'}\mathbf x_{\alpha'}\supset\mathbf F_{o\alpha'}(S_{\alpha'\alpha'}\mathbf x_{\alpha'})]
\supset.\,N_{o\alpha'}\mathbf x_{\alpha'}\supset\mathbf F_{o\alpha'}\mathbf x_{\alpha'}\).
Rule~V twice gives the conclusion.
\end{proof}
\end{annexitem}

\begin{annexitem}{annex:p02-w02-p02-040-church-27-29}{thm:p02-w02-p02-040-church-27-29}{Theorems 27, 28, 25 and 29}
\begin{proof}
\emph{27\(^o\)} \cite[\S6, pp.~64--65]{Church1940}.  Put
\(\mathbf a=[p\vee\sim p]\).  In 17\(^o\) substitute \(\mathbf a\) for
\(x_o\), \(\sim\mathbf a\) for \(y_o\) and
\(\lambda r[\sim[\mathbf a\supset\sim r]]\) for \(f_{oo}\), and reduce twice
by Rule~II:
\[
 \mathbf a=\sim\mathbf a\supset.\,\sim[\mathbf a\supset\sim\mathbf a]\supset\sim[\mathbf a\supset\sim\sim\mathbf a].
\]
The formula \(\sim[\mathbf a\supset\sim\mathbf a]\) is a tautology instance,
and (P), V with \((q\supset.\,r\supset s)\supset.\,r\supset.\,q\supset s\)
give \(\mathbf a=\sim\mathbf a\supset\sim[\mathbf a\supset\sim\sim\mathbf a]\).
The formula \(\mathbf a\supset\sim\sim\mathbf a\) is a tautology instance, and
(P), V with \(q\supset.\,(r\supset\sim q)\supset\sim r\) give
\(\mathbf a\neq\sim\mathbf a\).  Thirteen at type \(o\), with
\(\lambda t_o[\mathbf a\neq t_o]\) for the function variable and
\(\sim\mathbf a\) for the instance, gives
\((\exists t_o)[\mathbf a\neq t_o]\); thirteen again, with
\(\lambda s_o(\exists t_o)[s_o\neq t_o]\) and \(\mathbf a\), gives
\((\exists s_o)(\exists t_o)\,.\,s_o\neq t_o\), and Rule~I renames to 27\(^o\).

\emph{27\(^{\alpha\beta}\) from 27\(^\alpha\)} \cite[\S6, p.~65]{Church1940}.
By Rule~II, \(K_{\alpha\beta\alpha}z_\alpha x_\beta\) converts to \(z_\alpha\).
Hence, by Rule~III applied twice to the assumption,
\(z_\alpha\neq t_\alpha\derives K_{\alpha\beta\alpha}z_\alpha x_\beta\neq K_{\alpha\beta\alpha}t_\alpha x_\beta\).
In 17 at type \(\alpha\beta\) substitute \(K_{\alpha\beta\alpha}z_\alpha\),
\(K_{\alpha\beta\alpha}t_\alpha\) and
\(\lambda g_{\alpha\beta}[g_{\alpha\beta}x_\beta\neq K_{\alpha\beta\alpha}t_\alpha x_\beta]\)
and reduce; Rule~V twice gives
\(z_\alpha\neq t_\alpha,\ K_{\alpha\beta\alpha}z_\alpha=K_{\alpha\beta\alpha}t_\alpha
\derives K_{\alpha\beta\alpha}t_\alpha x_\beta\neq K_{\alpha\beta\alpha}t_\alpha x_\beta\).
The deduction theorem moves the second assumption into an antecedent;
16\(^\alpha\) gives
\(K_{\alpha\beta\alpha}t_\alpha x_\beta=K_{\alpha\beta\alpha}t_\alpha x_\beta\); and
(P), V with \((q\supset\sim r)\supset.\,r\supset\sim q\) give
\(z_\alpha\neq t_\alpha\derives K_{\alpha\beta\alpha}z_\alpha\neq K_{\alpha\beta\alpha}t_\alpha\).
Two applications of 13 at type \(\alpha\beta\), as in the case of \(o\), give
\(z_\alpha\neq t_\alpha\derives(\exists x_{\alpha\beta})(\exists y_{\alpha\beta})\,.\,x_{\alpha\beta}\neq y_{\alpha\beta}\).
No variable free in the assumption was substituted for or generalized
upon, so the deduction theorem gives
\(\derives z_\alpha\neq t_\alpha\supset\mathbf C\), where \(\mathbf C\) is the
existential just obtained, with no free \(z_\alpha,t_\alpha\).  Rule~E twice gives
\(\derives(\exists z_\alpha)(\exists t_\alpha)[z_\alpha\neq t_\alpha]\supset\mathbf C\),
and 27\(^\alpha\) with Rule~V gives 27\(^{\alpha\beta}\).

\emph{27\(^\alpha\) for every \(\alpha\).}  By induction on type symbols:
27\(^\iota\) is Axiom~7, 27\(^o\) was derived above, and 27\(^{\alpha\beta}\)
follows from 27\(^\alpha\).

\emph{28\(^\alpha\)} \cite[\S6, p.~65]{Church1940}.  Let \(z_\alpha,t_\alpha\) be
distinct from each other and from every variable of \(x_{\alpha'}\).
By Rule~II,
\(S_{\alpha'\alpha'}x_{\alpha'}(K_{\alpha\alpha\alpha}z_\alpha)t_\alpha\) converts to
\(K_{\alpha\alpha\alpha}z_\alpha(x_{\alpha'}(K_{\alpha\alpha\alpha}z_\alpha)t_\alpha)\)
and then to \(z_\alpha\), and \(0_{\alpha'}(K_{\alpha\alpha\alpha}z_\alpha)t_\alpha\)
converts to \(t_\alpha\).  So Rule~III gives
\(z_\alpha\neq t_\alpha\derives
S_{\alpha'\alpha'}x_{\alpha'}(K_{\alpha\alpha\alpha}z_\alpha)t_\alpha\neq0_{\alpha'}(K_{\alpha\alpha\alpha}z_\alpha)t_\alpha\).
In 17 at type \(\alpha'\) substitute \(S_{\alpha'\alpha'}x_{\alpha'}\),
\(0_{\alpha'}\) and
\(\lambda g_{\alpha'}[g_{\alpha'}(K_{\alpha\alpha\alpha}z_\alpha)t_\alpha\neq0_{\alpha'}(K_{\alpha\alpha\alpha}z_\alpha)t_\alpha]\),
and reduce; with the previous line, Rule~V gives
\(z_\alpha\neq t_\alpha,\ S_{\alpha'\alpha'}x_{\alpha'}=0_{\alpha'}\derives
0_{\alpha'}(K_{\alpha\alpha\alpha}z_\alpha)t_\alpha\neq0_{\alpha'}(K_{\alpha\alpha\alpha}z_\alpha)t_\alpha\).
As in the previous case, the deduction theorem, 16\(^\alpha\) and (P) give
\(z_\alpha\neq t_\alpha\derives S_{\alpha'\alpha'}x_{\alpha'}\neq0_{\alpha'}\);
the deduction theorem, Rule~E twice and 27\(^\alpha\) remove the assumption.

\emph{25\(^\alpha\).}  From 28\(^\alpha\) by (P), V with \(p\supset.\,q\supset p\).

\emph{29\(^\alpha\)} \cite[\S6, p.~65]{Church1940}.  Let \(\mathbf F\) be
\(\lambda x_{\alpha''}[N_{o\alpha'}(x_{\alpha''}S_{\alpha'\alpha'}0_{\alpha'})]\);
\(n_{\alpha''}\) is not free in it.  First, \(\mathbf F\,0_{\alpha''}\)
converts to \(N_{o\alpha'}(0_{\alpha''}S_{\alpha'\alpha'}0_{\alpha'})\) and then to
\(N_{o\alpha'}0_{\alpha'}\), because \(0_{\alpha''}\) is
\(\lambda f_{\alpha'\alpha'}\lambda x_{\alpha'}x_{\alpha'}\); so 22\(^\alpha\) and
Rule~III give \(\derives\mathbf F\,0_{\alpha''}\).  Second, on the
assumptions \(N_{o\alpha''}n_{\alpha''}\) and \(\mathbf F\,n_{\alpha''}\),
conversion gives \(N_{o\alpha'}(n_{\alpha''}S_{\alpha'\alpha'}0_{\alpha'})\);
23\(^\alpha\) with \(n_{\alpha''}S_{\alpha'\alpha'}0_{\alpha'}\) for
\(x_{\alpha'}\) and Rule~V give
\(N_{o\alpha'}(S_{\alpha'\alpha'}(n_{\alpha''}S_{\alpha'\alpha'}0_{\alpha'}))\).
Since \(S_{\alpha''\alpha''}n_{\alpha''}S_{\alpha'\alpha'}0_{\alpha'}\) converts
to \(S_{\alpha'\alpha'}(n_{\alpha''}S_{\alpha'\alpha'}0_{\alpha'})\), Rule~III
gives \(\mathbf F(S_{\alpha''\alpha''}n_{\alpha''})\).  The induction theorem
VIII at type \(\alpha''\), with no further assumptions, gives
\(N_{o\alpha''}n_{\alpha''}\supset\mathbf F\,n_{\alpha''}\), and Rule~II gives
29\(^\alpha\).
\end{proof}
\end{annexitem}

\paragraph{Conventions for the arithmetic derivations.}  In the next five
items, theorems 12--24 and 27--29 are used under assumptions through Step~1
of the deduction theorem, always after IV\('\) instances that substitute only
for their own variables.  Equality is handled by these instances:
from \(\Gamma\derives\mathbf A=\mathbf B\) follow
\(\Gamma\derives\mathbf B=\mathbf A\) (19),
\(\Gamma\derives\mathbf F\mathbf A=\mathbf F\mathbf B\) for any \(\mathbf F\) of
suitable type (18), and \(\Gamma\derives\mathbf G\mathbf B\) whenever
\(\Gamma\derives\mathbf G\mathbf A\) (17); with
\(\Gamma\derives\mathbf B=\mathbf C\) also \(\Gamma\derives\mathbf A=\mathbf C\)
(20).  Conversion (Rules I--III) is unrestricted under assumptions.  For a wff
\(\mathbf M\) of type \(\alpha''\) write
\(\mathbf M^\flat\) for \(\mathbf M\,S_{\alpha'\alpha'}0_{\alpha'}\).  Unfolding
\(S_{\alpha''\alpha''}\) and \(0_{\alpha''}\) by Rule~II gives the two
conversions
\[
 \text{(C1)}\quad(S_{\alpha''\alpha''}\mathbf M)^\flat\ \text{converts to}\
 S_{\alpha'\alpha'}(\mathbf M^\flat),\qquad
 \text{(C2)}\quad0_{\alpha''}^\flat\ \text{converts to}\ 0_{\alpha'}.
\]
A conjunction \(\mathbf A\mathbf B\) is split and formed with the tautologies
\(pq\supset p\), \(pq\supset q\) and \(p\supset.\,q\supset pq\) under (P).

\begin{annexitem}{annex:p02-w02-p02-040-church-26-30}{thm:p02-w02-p02-040-church-26-30}{Theorems 26 and 30 for the types \(\iota,\iota',\iota'',\ldots\)}
\begin{proof}
The route is the source's \cite[\S6, pp.~66--67]{Church1940}.

\emph{Clause 2.}  Assume \(\derives26^\alpha\) and \(\derives30^\alpha\).  Take
the assumptions \(N_{o\alpha''}x_{\alpha''}\), \(N_{o\alpha''}y_{\alpha''}\) and
\(S_{\alpha''\alpha''}x_{\alpha''}=S_{\alpha''\alpha''}y_{\alpha''}\).  By 18
with \(\lambda z_{\alpha''}(z_{\alpha''}S_{\alpha'\alpha'}0_{\alpha'})\) for the
function and Rule~II,
\((S_{\alpha''\alpha''}x_{\alpha''})^\flat=(S_{\alpha''\alpha''}y_{\alpha''})^\flat\),
which by (C1) converts to
\(S_{\alpha'\alpha'}(x_{\alpha''}^\flat)=S_{\alpha'\alpha'}(y_{\alpha''}^\flat)\).
By 29\(^\alpha\), \(N_{o\alpha'}(x_{\alpha''}^\flat)\) and
\(N_{o\alpha'}(y_{\alpha''}^\flat)\).  The instance of 26\(^\alpha\) with
\(x_{\alpha''}^\flat,y_{\alpha''}^\flat\) for \(x_{\alpha'},y_{\alpha'}\) and
Rule~V three times give \(x_{\alpha''}^\flat=y_{\alpha''}^\flat\).  The instance
of 30\(^\alpha\) with \(x_{\alpha''},y_{\alpha''}\) for \(m_{\alpha''},n_{\alpha''}\)
and Rule~V three times give \(x_{\alpha''}=y_{\alpha''}\).  No variable of an
assumption was substituted for or generalized upon, so the deduction theorem
three times gives 26\(^{\alpha'}\).

\emph{Clause 1.}  Assume \(\derives26^\alpha\).  Let \(\Phi(\mathbf m)\) be
\((n_{\alpha''})\,.\,N_{o\alpha''}n_{\alpha''}\supset.\,\mathbf m^\flat=n_{\alpha''}^\flat\supset\mathbf m=n_{\alpha''}\),
for wffs \(\mathbf m\) of type \(\alpha''\) in which \(n_{\alpha''}\) is not free.

\emph{(i) Zero against a successor.}  For every wff \(\mathbf k\) of type \(\alpha''\),
\(\derives0_{\alpha''}^\flat=(S_{\alpha''\alpha''}\mathbf k)^\flat\supset0_{\alpha''}=S_{\alpha''\alpha''}\mathbf k\)
and \(\derives(S_{\alpha''\alpha''}\mathbf k)^\flat=0_{\alpha''}^\flat\supset S_{\alpha''\alpha''}\mathbf k=0_{\alpha''}\).
For the first, by (C1) and (C2) the antecedent converts to
\(0_{\alpha'}=S_{\alpha'\alpha'}(\mathbf k^\flat)\); on it as assumption, 19 gives
\(S_{\alpha'\alpha'}(\mathbf k^\flat)=0_{\alpha'}\), while 28\(^\alpha\) with
\(\mathbf k^\flat\) for \(x_{\alpha'}\) gives
\(S_{\alpha'\alpha'}(\mathbf k^\flat)\neq0_{\alpha'}\); (P), V with
\(p\supset.\,\sim p\supset q\) give the consequent, and the deduction theorem
gives the implication.  The second is the same without 19.

\emph{(ii) The base \(\Phi(0_{\alpha''})\).}  Let \(\mathbf G\) be
\(\lambda n_{\alpha''}[0_{\alpha''}^\flat=n_{\alpha''}^\flat\supset0_{\alpha''}=n_{\alpha''}]\).
By 16, \(0_{\alpha''}=0_{\alpha''}\), and (P), V with \(q\supset.\,p\supset q\)
give \(\mathbf G\,0_{\alpha''}\) after Rule~III.  By (i),
\(\derives\mathbf G(S_{\alpha''\alpha''}n_{\alpha''})\) after Rule~III, so in
particular it follows from the assumptions \(N_{o\alpha''}n_{\alpha''}\) and
\(\mathbf G\,n_{\alpha''}\).  The induction theorem VIII gives
\(\derives N_{o\alpha''}n_{\alpha''}\supset\mathbf G\,n_{\alpha''}\); Rule~II and
VI\('\) give \(\Phi(0_{\alpha''})\).

\emph{(iii) The step.}  Take the assumptions \(N_{o\alpha''}m_{\alpha''}\) and
\(\Phi(m_{\alpha''})\); \(n_{\alpha''}\) is not free in them.  Let \(\mathbf H\) be
\(\lambda n_{\alpha''}[(S_{\alpha''\alpha''}m_{\alpha''})^\flat=n_{\alpha''}^\flat\supset S_{\alpha''\alpha''}m_{\alpha''}=n_{\alpha''}]\).
By (i), \(\derives\mathbf H\,0_{\alpha''}\) after Rule~III.  Add the assumptions
\(N_{o\alpha''}n_{\alpha''}\), \(\mathbf H\,n_{\alpha''}\) (unused) and
\((S_{\alpha''\alpha''}m_{\alpha''})^\flat=(S_{\alpha''\alpha''}n_{\alpha''})^\flat\).
By (C1) the last converts to
\(S_{\alpha'\alpha'}(m_{\alpha''}^\flat)=S_{\alpha'\alpha'}(n_{\alpha''}^\flat)\).
By 29\(^\alpha\), \(N_{o\alpha'}(m_{\alpha''}^\flat)\) and \(N_{o\alpha'}(n_{\alpha''}^\flat)\);
26\(^\alpha\) gives \(m_{\alpha''}^\flat=n_{\alpha''}^\flat\).  Twelve at type
\(\alpha''\) applied to \(\Phi(m_{\alpha''})\), with Rule~II, gives
\(N_{o\alpha''}n_{\alpha''}\supset.\,m_{\alpha''}^\flat=n_{\alpha''}^\flat\supset m_{\alpha''}=n_{\alpha''}\),
so Rule~V twice gives \(m_{\alpha''}=n_{\alpha''}\), and 18 with
\(S_{\alpha''\alpha''}\) gives \(S_{\alpha''\alpha''}m_{\alpha''}=S_{\alpha''\alpha''}n_{\alpha''}\).
The deduction theorem and Rule~III give
\(N_{o\alpha''}m_{\alpha''},\Phi(m_{\alpha''}),N_{o\alpha''}n_{\alpha''},\mathbf H\,n_{\alpha''}\derives\mathbf H(S_{\alpha''\alpha''}n_{\alpha''})\).
VIII on \(n_{\alpha''}\), with the first two as its assumptions, gives
\(N_{o\alpha''}m_{\alpha''},\Phi(m_{\alpha''})\derives N_{o\alpha''}n_{\alpha''}\supset\mathbf H\,n_{\alpha''}\),
and VI\('\) on \(n_{\alpha''}\) with Rule~II gives \(\Phi(S_{\alpha''\alpha''}m_{\alpha''})\).

\emph{(iv) Conclusion.}  VIII on \(m_{\alpha''}\) with
\(\lambda m_{\alpha''}\Phi(m_{\alpha''})\), (ii) and (iii) gives
\(\derives N_{o\alpha''}m_{\alpha''}\supset\Phi(m_{\alpha''})\).  Twelve at type
\(\alpha''\) gives
\(\Phi(m_{\alpha''})\supset.\,N_{o\alpha''}n_{\alpha''}\supset.\,m_{\alpha''}^\flat=n_{\alpha''}^\flat\supset m_{\alpha''}=n_{\alpha''}\),
and (P), V with \(p\supset q\supset.\,(q\supset r)\supset.\,p\supset r\) give 30\(^\alpha\).

\emph{The types \(\iota,\iota',\ldots\)}  Put \(\alpha_0=\iota\) and
\(\alpha_{k+1}=\alpha_k'\).  The formula 26\(^\iota\) is Axiom~8 literally.  If
\(\derives26^{\alpha_k}\), clause~1 gives \(\derives30^{\alpha_k}\) and clause~2
gives \(\derives26^{\alpha_{k+1}}\).  By induction on \(k\), 26\(^{\alpha_k}\) and
30\(^{\alpha_k}\) are theorems for every \(k\).  This discharges the
announcement of p.~64 that 26\(^{\iota'}\), 26\(^{\iota''}\), \ldots{} are
theorems.  The axioms used are 1--6\(^\gamma\), Axiom~8, and, through
28\(^{\alpha_k}\) and 27\(^\iota\), Axiom~7.
\end{proof}
\end{annexitem}

\begin{annexitem}{annex:p02-w02-p02-040-church-31-35}{thm:p02-w02-p02-040-church-31-35}{Theorems 31--35 on \(T_{\alpha''\alpha'}\)}
\begin{proof}
The route is the source's \cite[\S7, pp.~67--68]{Church1940}.  For a wff
\(\mathbf a\) of type \(\alpha'\) let \(\mathbf K(\mathbf z,\mathbf a)\) be
\(N_{o\alpha''}\mathbf z\,.\,\mathbf z^\flat=\mathbf a\), and
\(\mathbf G_{\mathbf a}\) the abstract \(\lambda x_{\alpha''}\mathbf K(x_{\alpha''},\mathbf a)\).
By the definition of \(T_{\alpha''\alpha'}\), Rule~II converts
\(T_{\alpha''\alpha'}\mathbf a\) to \(\iota_{\alpha''(o\alpha'')}\mathbf G_{\mathbf a}\).

\emph{33\(^\alpha\).}  Let \(\mathbf E\) be
\(\lambda x_{\alpha'}(\exists x_{\alpha''})\mathbf K(x_{\alpha''},x_{\alpha'})\);
the two variables differ in type.  Base: 16 at type \(\alpha'\) gives
\(0_{\alpha'}=0_{\alpha'}\), and Rule~III with (C2) turns the left side into
\(0_{\alpha''}^\flat\).  With 22 at type \(\alpha''\) and (P), V,
\(\mathbf K(0_{\alpha''},0_{\alpha'})\); 13 at type \(\alpha''\) with
\(\mathbf G_{0_{\alpha'}}\) and \(0_{\alpha''}\), and Rule~III, give
\(\mathbf E\,0_{\alpha'}\).  Step: on the assumption
\(\mathbf K(x_{\alpha''},x_{\alpha'})\), split it; 23 at type \(\alpha''\) gives
\(N_{o\alpha''}(S_{\alpha''\alpha''}x_{\alpha''})\); 18 with \(S_{\alpha'\alpha'}\)
gives \(S_{\alpha'\alpha'}(x_{\alpha''}^\flat)=S_{\alpha'\alpha'}x_{\alpha'}\), which
by (C1) is \((S_{\alpha''\alpha''}x_{\alpha''})^\flat=S_{\alpha'\alpha'}x_{\alpha'}\)
after Rule~III; so \(\mathbf K(S_{\alpha''\alpha''}x_{\alpha''},S_{\alpha'\alpha'}x_{\alpha'})\),
and 13 at type \(\alpha''\) gives
\((\exists x_{\alpha''})\mathbf K(x_{\alpha''},S_{\alpha'\alpha'}x_{\alpha'})\).
The deduction theorem and the derived rule E of
\cref{annex:p02-w02-p02-040-church-12-21} (generalizing on \(x_{\alpha''}\),
which is bound in the consequent) give
\(\derives(\exists x_{\alpha''})\mathbf K(x_{\alpha''},x_{\alpha'})\supset
(\exists x_{\alpha''})\mathbf K(x_{\alpha''},S_{\alpha'\alpha'}x_{\alpha'})\), whence
\(\mathbf E\,x_{\alpha'}\derives\mathbf E(S_{\alpha'\alpha'}x_{\alpha'})\) by
conversion and Rule~V.  VIII gives 33\(^\alpha\).  Only Axioms
1--6\(^\gamma\) were used.

\emph{Uniqueness, using 30\(^\alpha\).}  On the assumptions
\(\mathbf K(x_{\alpha''},x_{\alpha'})\) and \(\mathbf K(y_{\alpha''},x_{\alpha'})\),
splitting, 19 and 20 give \(x_{\alpha''}^\flat=y_{\alpha''}^\flat\), and the
instance of 30\(^\alpha\) at \(x_{\alpha''},y_{\alpha''}\) with Rule~V three times
gives \(x_{\alpha''}=y_{\alpha''}\).  The deduction theorem and VI\('\) on
\(y_{\alpha''}\), which is not free in the remaining assumption, give
\[
 \mathbf K(x_{\alpha''},x_{\alpha'})\derives(y_{\alpha''})\,.\,
 \mathbf K(y_{\alpha''},x_{\alpha'})\supset x_{\alpha''}=y_{\alpha''}.
\]

\emph{31\(^\alpha\) and 32\(^\alpha\).}  Substitute \(\mathbf G_{x_{\alpha'}}\) for
\(f_{o\alpha''}\) in 9\(^{\alpha''}\) and reduce: its free variable
\(x_{\alpha'}\) differs in type from the bound variables of 9\(^{\alpha''}\).
The result is
\(\mathbf K(x_{\alpha''},x_{\alpha'})\supset.\,(y_{\alpha''})[\mathbf K(y_{\alpha''},x_{\alpha'})\supset x_{\alpha''}=y_{\alpha''}]\supset
\mathbf G_{x_{\alpha'}}(\iota_{\alpha''(o\alpha'')}\mathbf G_{x_{\alpha'}})\).
With the uniqueness line, Rule~V twice, Rule~II and Rule~III on
\(T_{\alpha''\alpha'}x_{\alpha'}\),
\(\mathbf K(x_{\alpha''},x_{\alpha'})\derives\mathbf K(T_{\alpha''\alpha'}x_{\alpha'},x_{\alpha'})\).
The deduction theorem and rule E give
\(\derives(\exists x_{\alpha''})\mathbf K(x_{\alpha''},x_{\alpha'})\supset\mathbf K(T_{\alpha''\alpha'}x_{\alpha'},x_{\alpha'})\).
With 33\(^\alpha\), (P), V with
\(p\supset q\supset.\,[q\supset rs]\supset.\,p\supset r\) and
\(p\supset q\supset.\,[q\supset rs]\supset.\,p\supset s\) give 31\(^\alpha\) and
32\(^\alpha\).

\emph{34\(^\alpha\).}  Assume \(N_{o\alpha'}x_{\alpha'}\) and write \(\mathbf t\) for
\(T_{\alpha''\alpha'}x_{\alpha'}\).  By 31\(^\alpha\) and 23 at type \(\alpha''\),
\(N_{o\alpha''}(S_{\alpha''\alpha''}\mathbf t)\).  By 32\(^\alpha\),
\(\mathbf t^\flat=x_{\alpha'}\); 18 with \(S_{\alpha'\alpha'}\) and (C1) give
\((S_{\alpha''\alpha''}\mathbf t)^\flat=S_{\alpha'\alpha'}x_{\alpha'}\).  By
23\(^\alpha\), \(N_{o\alpha'}(S_{\alpha'\alpha'}x_{\alpha'})\); then 31\(^\alpha\)
and 32\(^\alpha\) at \(S_{\alpha'\alpha'}x_{\alpha'}\) give
\(N_{o\alpha''}(T_{\alpha''\alpha'}(S_{\alpha'\alpha'}x_{\alpha'}))\) and
\((T_{\alpha''\alpha'}(S_{\alpha'\alpha'}x_{\alpha'}))^\flat=S_{\alpha'\alpha'}x_{\alpha'}\).
By 19 and 20 the two \(\flat\)-values are equal, and 30\(^\alpha\) at
\(S_{\alpha''\alpha''}\mathbf t\) and \(T_{\alpha''\alpha'}(S_{\alpha'\alpha'}x_{\alpha'})\),
with 19, gives
\(T_{\alpha''\alpha'}(S_{\alpha'\alpha'}x_{\alpha'})=S_{\alpha''\alpha''}\mathbf t\).
The deduction theorem gives 34\(^\alpha\).

\emph{35\(^\alpha\).}  By 22\(^\alpha\), 31\(^\alpha\) and 32\(^\alpha\) at
\(0_{\alpha'}\): \(N_{o\alpha''}(T_{\alpha''\alpha'}0_{\alpha'})\) and
\((T_{\alpha''\alpha'}0_{\alpha'})^\flat=0_{\alpha'}\).  By 22 at type \(\alpha''\),
\(N_{o\alpha''}0_{\alpha''}\), and by 16 with (C2),
\(0_{\alpha''}^\flat=0_{\alpha'}\).  With 19 and 20, 30\(^\alpha\) at
\(T_{\alpha''\alpha'}0_{\alpha'}\) and \(0_{\alpha''}\) gives 35\(^\alpha\).
\end{proof}
\end{annexitem}

\begin{annexitem}{annex:p02-w02-p02-040-church-36-43}{thm:p02-w02-p02-040-church-36-43}{Theorems 36--43}
\begin{proof}
Write \(\mathbf i\) for \(I_{\alpha\alpha}\), \(\kappa_1\) for
\(K_{\alpha'\alpha'\alpha'}I_{\alpha'}\) and \(\kappa_0\) for
\(K_{\alpha'\alpha'\alpha'}0_{\alpha'}\).  Unfolding
\(\omega_{\alpha''\alpha'\alpha'}\), \(K\), \(I\) and \(0\) by Rule~II gives, for
all wffs \(\mathbf M,\mathbf N\) of type \(\alpha'\) (with \(h_{\alpha\alpha}\),
\(x_\alpha\) not free in them):
\[
\begin{aligned}
 \text{(\(\Pi\)1)}\quad&\langle\mathbf M,\mathbf N\rangle\kappa_1 0_{\alpha'}\ \text{converts to}\
 \lambda h_{\alpha\alpha}\lambda x_\alpha(\mathbf M h_{\alpha\alpha}(\mathbf N\mathbf i x_\alpha)),\\
 \text{(\(\Pi\)2)}\quad&\langle\mathbf M,\mathbf N\rangle\kappa_0 I_{\alpha'}\ \text{converts to}\
 \lambda h_{\alpha\alpha}\lambda x_\alpha(\mathbf M\mathbf i(\mathbf N h_{\alpha\alpha}x_\alpha)),
\end{aligned}
\]
since \(\kappa_1 0_{\alpha'}h_{\alpha\alpha}\) converts to \(h_{\alpha\alpha}\),
\(0_{\alpha'}h_{\alpha\alpha}\) to \(\mathbf i\), \(\kappa_0I_{\alpha'}h_{\alpha\alpha}\)
to \(\mathbf i\) and \(I_{\alpha'}h_{\alpha\alpha}\) to \(h_{\alpha\alpha}\).
Moreover \(S_{\alpha'\alpha'}\mathbf M\mathbf i\mathbf y\) converts to
\(\mathbf i(\mathbf M\mathbf i\mathbf y)\) and then to \(\mathbf M\mathbf i\mathbf y\),
and \(0_{\alpha'}\mathbf i\mathbf y\) converts to \(\mathbf y\).

\emph{36\(^\alpha\).}  \(\lambda f_{\alpha\alpha}\lambda x_\alpha(0_{\alpha'}f_{\alpha\alpha}x_\alpha)\)
converts to \(\lambda f_{\alpha\alpha}\lambda x_\alpha x_\alpha\), which is
\(0_{\alpha'}\); and both
\(\lambda f_{\alpha\alpha}\lambda x_\alpha(S_{\alpha'\alpha'}n_{\alpha'}f_{\alpha\alpha}x_\alpha)\)
and \(S_{\alpha'\alpha'}n_{\alpha'}\) convert to
\(\lambda f_{\alpha\alpha}\lambda x_\alpha(f_{\alpha\alpha}(n_{\alpha'}f_{\alpha\alpha}x_\alpha))\).
So 16 and conversion prove the formula
\(\lambda f_{\alpha\alpha}\lambda x_\alpha(\mathbf nf_{\alpha\alpha}x_\alpha)=\mathbf n\)
for \(\mathbf n=0_{\alpha'}\) and for \(\mathbf n=S_{\alpha'\alpha'}n_{\alpha'}\)
outright, and VIII gives 36\(^\alpha\).

\emph{37\(^\alpha\).}  On the assumption \(N_{o\alpha'}m_{\alpha'}\), induct on
\(n_{\alpha'}\) with \(\mathbf P(\mathbf n)\) the formula
\(\langle m_{\alpha'},\mathbf n\rangle\kappa_10_{\alpha'}=m_{\alpha'}\).  By
(\(\Pi\)1), \(\langle m_{\alpha'},0_{\alpha'}\rangle\kappa_10_{\alpha'}\) converts to
\(\lambda h_{\alpha\alpha}\lambda x_\alpha(m_{\alpha'}h_{\alpha\alpha}x_\alpha)\),
which equals \(m_{\alpha'}\) by 36\(^\alpha\); so \(\mathbf P(0_{\alpha'})\).
By (\(\Pi\)1) and the remark on \(S_{\alpha'\alpha'}\mathbf M\mathbf i\mathbf y\),
\(\langle m_{\alpha'},S_{\alpha'\alpha'}n_{\alpha'}\rangle\kappa_10_{\alpha'}\) and
\(\langle m_{\alpha'},n_{\alpha'}\rangle\kappa_10_{\alpha'}\) convert to the same
formula, so \(\mathbf P(S_{\alpha'\alpha'}n_{\alpha'})\) follows from
\(\mathbf P(n_{\alpha'})\) by conversion.  VIII (with \(n_{\alpha'}\) not free in
the assumption) and the deduction theorem give 37\(^\alpha\).

\emph{38\(^\alpha\).}  On the assumption \(N_{o\alpha'}n_{\alpha'}\), induct on
\(m_{\alpha'}\) with
\(\langle\mathbf m,n_{\alpha'}\rangle\kappa_0I_{\alpha'}=n_{\alpha'}\): by (\(\Pi\)2)
the case \(0_{\alpha'}\) converts to
\(\lambda h_{\alpha\alpha}\lambda x_\alpha(n_{\alpha'}h_{\alpha\alpha}x_\alpha)=n_{\alpha'}\),
which is 36\(^\alpha\), and the cases \(S_{\alpha'\alpha'}m_{\alpha'}\) and
\(m_{\alpha'}\) convert to the same formula.  VIII, the deduction theorem and
(P), V with \((p\supset.\,q\supset r)\supset.\,q\supset.\,p\supset r\) give 38\(^\alpha\).

\emph{39\(^\alpha\).}  Let \(\Psi\) be
\(\lambda p_{\alpha''}\langle S_{\alpha'\alpha'}(p_{\alpha''}\kappa_10_{\alpha'}),p_{\alpha''}\kappa_10_{\alpha'}\rangle\),
and for \(\mathbf n\) of type \(\alpha'''\) let \(\mathbf L(\mathbf n)\) be
\(\mathbf n\Psi\langle0_{\alpha'},0_{\alpha'}\rangle\kappa_10_{\alpha'}\) and
\(\mathbf R(\mathbf n)\) be
\((\mathbf nS_{\alpha''\alpha''}0_{\alpha''})^\flat\).  Induct on \(n_{\alpha'''}\)
(VIII at type \(\alpha''\)) with \(\mathbf L(\mathbf n)=\mathbf R(\mathbf n)\).
Base: \(\mathbf L(0_{\alpha'''})\) converts to
\(\langle0_{\alpha'},0_{\alpha'}\rangle\kappa_10_{\alpha'}\), which equals
\(0_{\alpha'}\) by 22\(^\alpha\) and 37\(^\alpha\); \(\mathbf R(0_{\alpha'''})\)
converts to \(0_{\alpha''}^\flat\) and then to \(0_{\alpha'}\).  Step: on the
assumptions \(N_{o\alpha'''}n_{\alpha'''}\) and
\(\mathbf L(n_{\alpha'''})=\mathbf R(n_{\alpha'''})\), write \(\mathbf u\) for
\(\mathbf L(n_{\alpha'''})\).  By 29 at type \(\alpha'\) and 29\(^\alpha\),
\(N_{o\alpha'}\mathbf R(n_{\alpha'''})\); so \(N_{o\alpha'}\mathbf u\) by 19 and
17, and \(N_{o\alpha'}(S_{\alpha'\alpha'}\mathbf u)\) by 23\(^\alpha\).  The
formula \(S_{\alpha'''\alpha'''}n_{\alpha'''}\Psi\langle0_{\alpha'},0_{\alpha'}\rangle\)
converts to \(\Psi(n_{\alpha'''}\Psi\langle0_{\alpha'},0_{\alpha'}\rangle)\) and
then to \(\langle S_{\alpha'\alpha'}\mathbf u,\mathbf u\rangle\), so
\(\mathbf L(S_{\alpha'''\alpha'''}n_{\alpha'''})\) converts to
\(\langle S_{\alpha'\alpha'}\mathbf u,\mathbf u\rangle\kappa_10_{\alpha'}\), which
equals \(S_{\alpha'\alpha'}\mathbf u\) by 37\(^\alpha\), hence
\(S_{\alpha'\alpha'}\mathbf R(n_{\alpha'''})\) by 18.  And
\(\mathbf R(S_{\alpha'''\alpha'''}n_{\alpha'''})\) converts to
\((S_{\alpha''\alpha''}(n_{\alpha'''}S_{\alpha''\alpha''}0_{\alpha''}))^\flat\) and,
by (C1), to \(S_{\alpha'\alpha'}\mathbf R(n_{\alpha'''})\).  By 19 and 20 the
step holds.  VIII gives 39\(^\alpha\).

\emph{40\(^\alpha\).}  With \(\mathbf u\) as in the step just given and on the
assumption \(N_{o\alpha'''}n_{\alpha'''}\), the formula
\(P_{\alpha'\alpha'''}(S_{\alpha'''\alpha'''}n_{\alpha'''})\) converts to
\(\langle S_{\alpha'\alpha'}\mathbf u,\mathbf u\rangle\kappa_0I_{\alpha'}\).  By
39\(^\alpha\), \(\mathbf u=\mathbf R(n_{\alpha'''})\); as there,
\(N_{o\alpha'}\mathbf u\) and \(N_{o\alpha'}(S_{\alpha'\alpha'}\mathbf u)\); 38\(^\alpha\)
gives \(\langle S_{\alpha'\alpha'}\mathbf u,\mathbf u\rangle\kappa_0I_{\alpha'}=\mathbf u\),
and 20 with the deduction theorem gives 40\(^\alpha\).

\emph{41\(^\alpha\).}  \(P_{\alpha'\alpha'''}0_{\alpha'''}\) converts to
\(\langle0_{\alpha'},0_{\alpha'}\rangle\kappa_0I_{\alpha'}\), by (\(\Pi\)2) to
\(\lambda h_{\alpha\alpha}\lambda x_\alpha(0_{\alpha'}\mathbf i(0_{\alpha'}h_{\alpha\alpha}x_\alpha))\),
and then to \(\lambda h_{\alpha\alpha}\lambda x_\alpha x_\alpha\), an alphabetic
variant of \(0_{\alpha'}\).  So 16 and conversion give 41\(^\alpha\).

\emph{42\(^\alpha\).}  Assume \(\derives30^\alpha\) and \(\derives30^{\alpha'}\), so
that 31--35 hold at \(\alpha\) and at \(\alpha'\).  \(P_{\alpha'\alpha'}0_{\alpha'}\)
converts to \(P_{\alpha'\alpha'''}(T_{\alpha'''\alpha''}(T_{\alpha''\alpha'}0_{\alpha'}))\).
By 35\(^\alpha\) and 18 with \(T_{\alpha'''\alpha''}\),
\(T_{\alpha'''\alpha''}(T_{\alpha''\alpha'}0_{\alpha'})=T_{\alpha'''\alpha''}0_{\alpha''}\),
and by 35\(^{\alpha'}\) and 20 this equals \(0_{\alpha'''}\).  Then 18 with
\(P_{\alpha'\alpha'''}\), 41\(^\alpha\) and 20 give 42\(^\alpha\).

\emph{43\(^\alpha\).}  Assume \(N_{o\alpha'}n_{\alpha'}\) and write \(\mathbf t\) for
\(T_{\alpha''\alpha'}n_{\alpha'}\) and \(\mathbf t'\) for \(T_{\alpha'''\alpha''}\mathbf t\).
By 31\(^\alpha\), \(N_{o\alpha''}\mathbf t\), and by 31\(^{\alpha'}\),
\(N_{o\alpha'''}\mathbf t'\).  By 34\(^\alpha\),
\(T_{\alpha''\alpha'}(S_{\alpha'\alpha'}n_{\alpha'})=S_{\alpha''\alpha''}\mathbf t\);
18 with \(T_{\alpha'''\alpha''}\), 34\(^{\alpha'}\) at \(\mathbf t\) and 20 give
\(T_{\alpha'''\alpha''}(T_{\alpha''\alpha'}(S_{\alpha'\alpha'}n_{\alpha'}))=S_{\alpha'''\alpha'''}\mathbf t'\).
Since \(P_{\alpha'\alpha'}(S_{\alpha'\alpha'}n_{\alpha'})\) converts to
\(P_{\alpha'\alpha'''}(T_{\alpha'''\alpha''}(T_{\alpha''\alpha'}(S_{\alpha'\alpha'}n_{\alpha'})))\),
18 with \(P_{\alpha'\alpha'''}\) and 40\(^\alpha\) at \(\mathbf t'\) give
\(P_{\alpha'\alpha'}(S_{\alpha'\alpha'}n_{\alpha'})=(\mathbf t'S_{\alpha''\alpha''}0_{\alpha''})^\flat\).
By 32\(^{\alpha'}\) at \(\mathbf t\),
\(\mathbf t'S_{\alpha''\alpha''}0_{\alpha''}=\mathbf t\); 18 with
\(\lambda z_{\alpha''}(z_{\alpha''}^\flat)\) gives
\((\mathbf t'S_{\alpha''\alpha''}0_{\alpha''})^\flat=\mathbf t^\flat\); and by
32\(^\alpha\), \(\mathbf t^\flat=n_{\alpha'}\).  Theorem 20 twice and the deduction
theorem give 43\(^\alpha\).
\end{proof}
\end{annexitem}

\begin{annexitem}{annex:p02-w02-p02-040-church-recursion}{thm:p02-w02-p02-040-church-recursion}{The recursion equations for Church's formula}
\begin{proof}
Notation as in \cref{annex:p02-w02-p02-040-church-36-43}.  Let \(\Psi_{\mathbf B}\)
be
\(\lambda y_{\alpha''}\langle S_{\alpha'\alpha'}(y_{\alpha''}\kappa_10_{\alpha'}),
\mathbf B_{\alpha'\alpha'\alpha'}(y_{\alpha''}\kappa_10_{\alpha'})(y_{\alpha''}\kappa_0I_{\alpha'})\rangle\)
and, for \(\mathbf n\) of type \(\alpha'''\), \(\mathbf W(\mathbf n)\) the formula
\(\mathbf n\Psi_{\mathbf B}\langle0_{\alpha'},\mathbf A_{\alpha'}\rangle\), with
first and second projections \(\mathbf u(\mathbf n)=\mathbf W(\mathbf n)\kappa_10_{\alpha'}\)
and \(\mathbf v(\mathbf n)=\mathbf W(\mathbf n)\kappa_0I_{\alpha'}\).  Then
\(\mathbf F_{\alpha'\alpha'}\mathbf a\) converts to
\(\mathbf v(T_{\alpha'''\alpha''}(T_{\alpha''\alpha'}\mathbf a))\) for every
\(\mathbf a\) of type \(\alpha'\).

\emph{Invariant.}  We show
\(\derives N_{o\alpha'''}n_{\alpha'''}\supset.\,\mathbf u(n_{\alpha'''})=\mathbf R(n_{\alpha'''})\,.\,N_{o\alpha'}\mathbf v(n_{\alpha'''})\),
with \(\mathbf R\) as in 39, by VIII on \(n_{\alpha'''}\).  Base:
\(\mathbf W(0_{\alpha'''})\) converts to \(\langle0_{\alpha'},\mathbf A_{\alpha'}\rangle\).
By (H), \(N_{o\alpha'}\mathbf A_{\alpha'}\); with 22\(^\alpha\), 37\(^\alpha\) and
38\(^\alpha\) give \(\mathbf u(0_{\alpha'''})=0_{\alpha'}\), which is
\(\mathbf R(0_{\alpha'''})\) up to conversion, and
\(\mathbf v(0_{\alpha'''})=\mathbf A_{\alpha'}\), whence
\(N_{o\alpha'}\mathbf v(0_{\alpha'''})\) by 19 and 17.  Step: on the assumptions
\(N_{o\alpha'''}n_{\alpha'''}\) and the invariant at \(n_{\alpha'''}\), write
\(\mathbf u,\mathbf v\) for \(\mathbf u(n_{\alpha'''}),\mathbf v(n_{\alpha'''})\).
As in the proof of 39, \(N_{o\alpha'}\mathbf u\) and
\(N_{o\alpha'}(S_{\alpha'\alpha'}\mathbf u)\); \(N_{o\alpha'}\mathbf v\) is assumed;
and (H) at \(\mathbf u,\mathbf v\) gives \(N_{o\alpha'}(\mathbf B_{\alpha'\alpha'\alpha'}\mathbf u\mathbf v)\).
\(\mathbf W(S_{\alpha'''\alpha'''}n_{\alpha'''})\) converts to
\(\langle S_{\alpha'\alpha'}\mathbf u,\mathbf B_{\alpha'\alpha'\alpha'}\mathbf u\mathbf v\rangle\),
so 37\(^\alpha\) and 38\(^\alpha\) give
\(\mathbf u(S_{\alpha'''\alpha'''}n_{\alpha'''})=S_{\alpha'\alpha'}\mathbf u\) and
\(\mathbf v(S_{\alpha'''\alpha'''}n_{\alpha'''})=\mathbf B_{\alpha'\alpha'\alpha'}\mathbf u\mathbf v\).
The first equals \(\mathbf R(S_{\alpha'''\alpha'''}n_{\alpha'''})\) exactly as in
the step of 39, and the second is a number by 19 and 17.

\emph{First equation.}  By 35\(^\alpha\), 35\(^{\alpha'}\), 18 and 20,
\(T_{\alpha'''\alpha''}(T_{\alpha''\alpha'}0_{\alpha'})=0_{\alpha'''}\); 18 with
\(\lambda n_{\alpha'''}\mathbf v(n_{\alpha'''})\) (the variable \(n_{\alpha'''}\)
is not free in \(\mathbf A_{\alpha'},\mathbf B_{\alpha'\alpha'\alpha'}\)) and
conversion give \(\mathbf F_{\alpha'\alpha'}0_{\alpha'}=\mathbf v(0_{\alpha'''})\),
which equals \(\mathbf A_{\alpha'}\) by the base case.

\emph{Second equation.}  Assume \(N_{o\alpha'}\mathbf x_{\alpha'}\), and let
\(\mathbf t'\) be \(T_{\alpha'''\alpha''}(T_{\alpha''\alpha'}\mathbf x_{\alpha'})\).  As
in the proof of 43, \(N_{o\alpha'''}\mathbf t'\),
\(T_{\alpha'''\alpha''}(T_{\alpha''\alpha'}(S_{\alpha'\alpha'}\mathbf x_{\alpha'}))=S_{\alpha'''\alpha'''}\mathbf t'\),
and \(\mathbf R(\mathbf t')=\mathbf x_{\alpha'}\).  By 18 and conversion,
\(\mathbf F_{\alpha'\alpha'}(S_{\alpha'\alpha'}\mathbf x_{\alpha'})=\mathbf v(S_{\alpha'''\alpha'''}\mathbf t')\).
The invariant at \(\mathbf t'\) (an instance by IV\('\) followed by Rule~V) and
the step computation give
\(\mathbf v(S_{\alpha'''\alpha'''}\mathbf t')=\mathbf B_{\alpha'\alpha'\alpha'}\mathbf u(\mathbf t')\mathbf v(\mathbf t')\),
with \(\mathbf u(\mathbf t')=\mathbf R(\mathbf t')=\mathbf x_{\alpha'}\).  Finally
\(\mathbf F_{\alpha'\alpha'}\mathbf x_{\alpha'}\) converts to \(\mathbf v(\mathbf t')\).
Applying 18 twice (to the two arguments of \(\mathbf B_{\alpha'\alpha'\alpha'}\))
and 20 gives
\(\mathbf F_{\alpha'\alpha'}(S_{\alpha'\alpha'}\mathbf x_{\alpha'})=
\mathbf B_{\alpha'\alpha'\alpha'}\mathbf x_{\alpha'}(\mathbf F_{\alpha'\alpha'}\mathbf x_{\alpha'})\),
and the deduction theorem finishes.  The assumption variable
\(\mathbf x_{\alpha'}\) was never substituted for or generalized upon.

\emph{Why a hypothesis is needed.}  Work in \(\mathcal S(\N,c)\) of
\cref{def:p02-w02-p02-040-church-standard} with \(\alpha=\iota\); every
theorem is valid there, and since \(\N\) is infinite the numerals of each type
\(\iota,\iota',\ldots\) are pairwise distinct (proof of
\cref{thm:p02-w02-p02-040-church-soundness}), so
\(T_{\iota''\iota'}\) and \(T_{\iota'''\iota''}\) send each numeral to the
corresponding numeral.  Let \(a_0\) be \(c_\iota\) applied to the constant function with value \(\mathrm T\), the value
of the closed wff \(\iota_{\iota(o\iota)}(\lambda z_\iota[z_\iota=z_\iota])\), and
let \(\mathbf A_{\iota'}\) be
\(\lambda h_{\iota\iota}\lambda x_\iota(\iota_{\iota(o\iota)}(\lambda z_\iota[z_\iota=z_\iota]))\),
which denotes the functional \(\kappa\) with \(\kappa(h)(x)=a_0\).  Let
\(\mathbf B_{\iota'\iota'\iota'}\) be \(K_{\iota'\iota'\iota'}\).  By (\(\Pi\)1), the
counter of \(\langle0_{\iota'},\mathbf A_{\iota'}\rangle\) denotes
\(h\mapsto x\mapsto\kappa(\mathrm{id})(x)=a_0\), that is \(\kappa\) again, not
\(\bar0\).  So \(\mathbf W(1_{\iota'''})\) denotes the pair
\(\langle S(\kappa),\kappa\rangle\), and by (\(\Pi\)2)
\(\mathbf F_{\iota'\iota'}(S_{\iota'\iota'}0_{\iota'})\) denotes
\(h\mapsto x\mapsto S(\kappa)(\mathrm{id})(\kappa(h)(x))=a_0\), which is
\(\kappa\).  But \(\mathbf B_{\iota'\iota'\iota'}0_{\iota'}(\mathbf F_{\iota'\iota'}0_{\iota'})\)
denotes \(\bar0\), and \(\kappa\neq\bar0\) because \(\bar0(h)(x)=x\neq a_0\) for
any \(x\neq a_0\).  So the second equation fails at \(0_{\iota'}\), where
\(N_{o\iota'}0_{\iota'}\) holds, and that instance is not a theorem.
\end{proof}
\end{annexitem}

\begin{annexitem}{annex:p02-w02-p02-040-church-rosser}{thm:p02-w02-p02-040-church-rosser}{Rosser's addition and multiplication}
\begin{proof}
Unfold \(\mathfrak s\) and \(\mathfrak B\) by Rule~II.
\(0_{\alpha'}+n_{\alpha'}\) converts to
\(\lambda f_{\alpha\alpha}\lambda x_\alpha(n_{\alpha'}f_{\alpha\alpha}x_\alpha)\),
which equals \(n_{\alpha'}\) under \(N_{o\alpha'}n_{\alpha'}\) by 36\(^\alpha\).
Both \(S_{\alpha'\alpha'}m_{\alpha'}+n_{\alpha'}\) and
\(S_{\alpha'\alpha'}(m_{\alpha'}+n_{\alpha'})\) convert to
\(\lambda f_{\alpha\alpha}\lambda x_\alpha(f_{\alpha\alpha}(m_{\alpha'}f_{\alpha\alpha}(n_{\alpha'}f_{\alpha\alpha}x_\alpha)))\),
so 16 and conversion give the second equation.
\(0_{\alpha'}\times n_{\alpha'}\) converts to
\(\lambda f_{\alpha\alpha}(0_{\alpha'}(n_{\alpha'}f_{\alpha\alpha}))\) and then to
\(\lambda f_{\alpha\alpha}\lambda x_\alpha x_\alpha\), which is \(0_{\alpha'}\).
Both \(S_{\alpha'\alpha'}m_{\alpha'}\times n_{\alpha'}\) and
\(n_{\alpha'}+(m_{\alpha'}\times n_{\alpha'})\) convert to
\(\lambda f_{\alpha\alpha}\lambda x_\alpha(n_{\alpha'}f_{\alpha\alpha}(m_{\alpha'}(n_{\alpha'}f_{\alpha\alpha})x_\alpha))\),
which gives the fourth equation by 16 and conversion.

Closure under \(+\): on the assumption \(N_{o\alpha'}n_{\alpha'}\), induct on
\(m_{\alpha'}\) with \(N_{o\alpha'}(\mathbf m+n_{\alpha'})\).  The base follows
from \(0_{\alpha'}+n_{\alpha'}=n_{\alpha'}\) by 19 and 17; in the step,
\(N_{o\alpha'}(m_{\alpha'}+n_{\alpha'})\) gives
\(N_{o\alpha'}(S_{\alpha'\alpha'}(m_{\alpha'}+n_{\alpha'}))\) by 23\(^\alpha\), and
conversion gives \(N_{o\alpha'}(S_{\alpha'\alpha'}m_{\alpha'}+n_{\alpha'})\).
Closure under \(\times\): same induction; the base is 22\(^\alpha\) after
conversion, and in the step closure under \(+\) at \(n_{\alpha'}\) and
\(m_{\alpha'}\times n_{\alpha'}\) gives \(N_{o\alpha'}(n_{\alpha'}+(m_{\alpha'}\times n_{\alpha'}))\),
which converts to \(N_{o\alpha'}(S_{\alpha'\alpha'}m_{\alpha'}\times n_{\alpha'})\).
VIII, the deduction theorem and the permutation tautology give the two
closure theorems.
\end{proof}
\end{annexitem}

\begin{annexitem}{annex:p02-w02-p02-040-church-soundness}{thm:p02-w02-p02-040-church-soundness}{Soundness, consistency and independence}
\begin{proof}
Fix \(\mathcal S(D,c)\).

\emph{Coincidence.}  If \(s\) and \(s'\) agree on the free variables of
\(\mathbf M\), then \(\sem{\mathbf M}_s=\sem{\mathbf M}_{s'}\).  Induction on
\(\mathbf M\); in the abstraction case, \(s[\mathbf x\mapsto a]\) and
\(s'[\mathbf x\mapsto a]\) agree on the free variables of the body.

\emph{Substitution.}  If the bound variables of \(\mathbf M\) are distinct
from \(\mathbf x\) and from the free variables of \(\mathbf N\), and
\(\mathbf M'\) is the result of substituting \(\mathbf N\) for
\(\mathbf x\) throughout \(\mathbf M\), then
\(\sem{\mathbf M'}_s=\sem{\mathbf M}_{s[\mathbf x\mapsto\sem{\mathbf N}_s]}\).
Induction on \(\mathbf M\).  The variable, constant and application cases
are immediate.  For \(\mathbf M=(\lambda\mathbf y\mathbf M_1)\) the hypothesis
gives \(\mathbf y\neq\mathbf x\) and \(\mathbf y\) not free in \(\mathbf N\);
for every \(a\),
\(\sem{\mathbf M_1'}_{s[\mathbf y\mapsto a]}
=\sem{\mathbf M_1}_{s[\mathbf y\mapsto a][\mathbf x\mapsto\sem{\mathbf N}_{s[\mathbf y\mapsto a]}]}
=\sem{\mathbf M_1}_{s[\mathbf x\mapsto\sem{\mathbf N}_s][\mathbf y\mapsto a]}\)
by the induction hypothesis, coincidence and \(\mathbf x\neq\mathbf y\).

\emph{Renaming.}  If \(\mathbf y\) does not occur in \(\mathbf M\) and
\(\mathbf M^{\mathbf y}\) results from substituting \(\mathbf y\) for
\(\mathbf x\) throughout \(\mathbf M\), binder occurrences included, then
\(\sem{\mathbf M^{\mathbf y}}_s=\sem{\mathbf M}_{s[\mathbf x\mapsto s(\mathbf y)]}\).
Induction on \(\mathbf M\); for \(\mathbf M=(\lambda\mathbf x\mathbf M_1)\),
both sides are \(a\mapsto\sem{\mathbf M_1}_{s[\mathbf x\mapsto a]}\) by the
induction hypothesis and coincidence, since \(\mathbf y\) does not occur in
\(\mathbf M_1\).  When \(\mathbf x\) is not free in \(\mathbf M\), coincidence
gives \(\sem{\mathbf M^{\mathbf y}}_s=\sem{\mathbf M}_s\).

\emph{Replacement.}  If \(\sem{\mathbf M}_s=\sem{\mathbf M'}_s\) for every
\(s\), then replacing one part \(\mathbf M\) of \(\mathbf P\) by
\(\mathbf M'\) gives \(\mathbf P'\) with \(\sem{\mathbf P}_s=\sem{\mathbf P'}_s\)
for every \(s\), by induction on \(\mathbf P\).  Parts exclude the variable
after \(\lambda\), so the abstraction case only replaces inside the body.

\emph{Rules.}  Rule~I preserves denotation by renaming and replacement.
Rules II and III preserve it because, under their proviso,
\(\sem{(\lambda\mathbf x_\beta\mathbf M)\mathbf N}_s=
\sem{\mathbf M}_{s[\mathbf x_\beta\mapsto\sem{\mathbf N}_s]}=\sem{\mathbf M'}_s\)
by substitution, and by replacement.  For Rule~IV, let
\(\mathbf F\mathbf x\) be valid with \(\mathbf x\) not free in \(\mathbf F\),
and let \(s\) be given; with \(s'=s[\mathbf x\mapsto\sem{\mathbf A}_s]\),
coincidence gives
\(\sem{\mathbf F\mathbf A}_s=\sem{\mathbf F}_{s'}(s'(\mathbf x))=\sem{\mathbf F\mathbf x}_{s'}=\mathrm T\).
Rule~V preserves validity by the meaning of \(\vee\) and \(\sim\).  For
Rule~VI, for every \(a\),
\(\sem{\mathbf F}_s(a)=\sem{\mathbf F\mathbf x}_{s[\mathbf x\mapsto a]}=\mathrm T\),
so \(\sem{\Pi\mathbf F}_s=\mathrm T\).

\emph{Equality.}  \(\sem{Q_{o\alpha\alpha}\mathbf A\mathbf B}_s=\mathrm T\) iff
every \(g\in\mathcal D_{o\alpha}\) with \(g(a)=\mathrm T\) has
\(g(b)=\mathrm T\), where \(a,b\) are the two denotations.  If \(a=b\) this
holds.  If it holds, take \(g\) to be the characteristic function of
\(\{a\}\), a member of \(\mathcal D_{o\alpha}\) because that set contains
every function; then \(g(b)=\mathrm T\), so \(b=a\).

\emph{Axioms 1--6, 9--11 and 7.}  The value of a wff built by \(\sim\) and
\(\vee\) from wffs of type \(o\) is computed by the Boolean truth tables, so
1--4 are valid.  Axiom~5: if \(\Pi f\) is true then \(f(a)=\mathrm T\) for
every \(a\), in particular \(a=s(x_\alpha)\).  Axiom~6: if
\(s(p)=\mathrm F\) and \(s(p)\vee f(a)\) is true for every \(a\), then
\(\Pi f\) is true.  Axiom~9 and Axiom~11: if \(f(a)=\mathrm T\) for some
\(a\), then \(f(c_\alpha(f))=\mathrm T\) by the choice of \(c_\alpha\).
Axiom~10: by the equality clause the antecedent says
\(f(a)=g(a)\) for every \(a\in\mathcal D_\beta\), so \(f=g\) as functions,
and the consequent is true.  Axiom~7 is true iff \(D\) has two distinct
members, by the equality clause.

\emph{Numerals.}  Write \(\varphi^n\) for the \(n\)-fold composite of
\(\varphi\in\mathcal D_{\alpha\alpha}\), with \(\varphi^0\) the identity, and
\(\bar n\) for the function \(\varphi\mapsto\varphi^n\) in
\(\mathcal D_{\alpha'}\).  By the recursion,
\(\sem{n_{\alpha'}}_s=\bar n\) and
\(\sem{S_{\alpha'\alpha'}}_s(m)=(\varphi\mapsto\varphi\circ m(\varphi))\), so
\(\sem{S_{\alpha'\alpha'}}_s(\bar n)=\overline{n+1}\).  Unfolding
\(N_{o\alpha'}\), \(\sem{N_{o\alpha'}}_s(m)=\mathrm T\) iff \(m\) belongs to every
subset of \(\mathcal D_{\alpha'}\) containing \(\bar0\) and closed under
\(\sem{S_{\alpha'\alpha'}}\) (all subsets are available, as characteristic
functions).  The set \(\{\bar n:n\in\N\}\) is such a subset and, by
induction on \(n\), is contained in every such subset; so
\(\sem{N_{o\alpha'}}_s(m)=\mathrm T\) iff \(m=\bar n\) for some \(n\).
Finally, if \(u\neq v\) in \(\mathcal D_\alpha\) and \(\kappa_u\) is the
constant function with value \(u\), then for every \(m\in\mathcal D_{\alpha'}\),
\(\sem{S_{\alpha'\alpha'}}_s(m)(\kappa_u)(v)=\kappa_u(m(\kappa_u)(v))=u\neq v=\bar0(\kappa_u)(v)\),
so \(\sem{S_{\alpha'\alpha'}}_s(m)\neq\bar0\).

\emph{Axiom~8 and 26\(^\alpha\).}  By the numeral clause, 26\(^\alpha\) is valid
iff \(\sem{S_{\alpha'\alpha'}}\) is injective on \(\{\bar n:n\in\N\}\), that
is, iff \(\overline{n+1}=\overline{m+1}\) implies \(\bar n=\bar m\); Axiom~8 is
26\(^\iota\).  If \(\mathcal D_\alpha\) has one member, \(\mathcal D_{\alpha'}\)
has one member and the condition holds.  If \(\mathcal D_\alpha\) is infinite,
choose distinct \(a_0,a_1,\ldots\) in it and let \(\varphi(a_i)=a_{i+1}\),
\(\varphi\) the identity elsewhere; then \(\bar n(\varphi)(a_0)=a_n\), so the
\(\bar n\) are pairwise distinct and the condition holds.  If
\(\mathcal D_\alpha\) is finite with at least two members, then
\(\mathcal D_{\alpha'}\) is finite, so \(\bar n=\bar m\) for some \(n<m\).
Suppose the condition held, and take the least such \(n\).  If \(n>0\), the
condition gives \(\overline{n-1}=\overline{m-1}\), contradicting minimality;
if \(n=0\), then \(\bar0=\bar m=\sem{S_{\alpha'\alpha'}}(\overline{m-1})\),
contradicting the last numeral clause.  So the condition fails, and
26\(^\alpha\) is not valid.

\emph{Conclusions.}  Clause~1 is now proved.  For clause~2, take
\(D=\N\): every axiom is valid, and the rules preserve validity, so every
theorem is valid; the variable \(p\) is false under an assignment with
\(s(p)=\mathrm F\), so it is not a theorem.  For clause~3, in
\(\mathcal S(D,c)\) with \(D\) a singleton every axiom other than 7 is valid
and 7 is not, and with \(D\) of exactly two members every axiom other than 8
is valid and 8 is not; validity is preserved by the rules, so neither axiom
is derivable from the others.  For clause~4, if \(\iota\) does not occur in
\(\alpha\), then by induction on \(\alpha\) the set \(\mathcal D_\alpha\) is
finite with at least two members, since \(\mathcal D_o\) has two members and
\(\mathcal D_{\gamma\delta}\) is finite with at least two members when
\(\mathcal D_\gamma\) and \(\mathcal D_\delta\) are finite and
\(\mathcal D_\gamma\) has at least two members.  Hence 26\(^\alpha\) is not
valid in \(\mathcal S(\N,c)\), where every theorem is valid.
\end{proof}
\end{annexitem}

\begin{annexitem}{annex:p02-w02-p02-040-church-9-from-11}{thm:p02-w02-p02-040-church-9-from-11}{Descriptions from choice}
\begin{proof}
The formula \((p\supset r)\supset.\,p\supset.\,q\supset r\) is a tautology in
\(\sim,\vee\).  By \cref{thm:p02-w02-p02-040-church-prop-complete}, which uses
only Axioms~1--4 and Rules I, IV\('\), V, its instance with
\(f_{o\alpha}x_\alpha\), \((y_\alpha)[f_{o\alpha}y_\alpha\supset x_\alpha=y_\alpha]\)
and \(f_{o\alpha}(\iota_{\alpha(o\alpha)}f_{o\alpha})\) for \(p,q,r\) is a theorem.
Its antecedent is Axiom 11\(^\alpha\), and Rule~V gives 9\(^\alpha\).
\end{proof}
\end{annexitem}

\begin{annexitem}{annex:p02-w02-p02-040-church-selection}{thm:p02-w02-p02-040-church-selection}{The defined selection constants satisfy 9; independence of 9}
\begin{proof}
\emph{Clause 1.}  By induction on type symbols.  9\(^o\) and 9\(^\iota\) are
axioms of the reduced system.  Let \(\gamma=\alpha\beta\) and assume 9\(^\alpha\)
is a theorem of the reduced system, with \(\iota_{\alpha(o\alpha)}\) primitive or
defined.  Write \(h\) for \(h_{o(\alpha\beta)}\), \(f,g\) for variables of type
\(\alpha\beta\), and let \(\mathbf G\) be
\(\lambda x_\beta(\text{\rotatebox[origin=c]{180}{\(\iota\)}}y_\alpha)(\exists g)\,.\,hg\,.\,y_\alpha=gx_\beta\),
to which \(\iota_{\alpha\beta(o(\alpha\beta))}h\) converts by Rule~II (after
renaming the bound \(f_{\alpha\beta}\) of the schema to \(g\)).  For the
variable \(x_\beta\) let \(\mathbf E\) be
\(\lambda y_\alpha(\exists g)\,.\,hg\,.\,y_\alpha=gx_\beta\), of type \(o\alpha\);
the body of \(\mathbf G\) is \(\iota_{\alpha(o\alpha)}\mathbf E\).  Take the
assumptions (a) \(hf\) and (b) \((g)[hg\supset f=g]\).
\begin{enumerate}
 \item \(\mathbf E(fx_\beta)\): by 16, \(fx_\beta=fx_\beta\); with (a) and
   (P), V, \(hf\,.\,fx_\beta=fx_\beta\); 13 at type \(\alpha\beta\) with
   \(\lambda g[hg\,.\,fx_\beta=gx_\beta]\) and \(f\) gives
   \((\exists g)[hg\,.\,fx_\beta=gx_\beta]\), which is \(\mathbf E(fx_\beta)\)
   after Rule~III.
 \item \((y_\alpha)[\mathbf Ey_\alpha\supset fx_\beta=y_\alpha]\) from (b): on the
   further assumption \(hg\,.\,y_\alpha=gx_\beta\), split it; 12 at type
   \(\alpha\beta\) applied to (b) and Rule~V give \(f=g\); 17 at type
   \(\alpha\beta\) with \(\lambda k_{\alpha\beta}[fx_\beta=kx_\beta]\), together
   with 16, gives \(fx_\beta=gx_\beta\); 19 and 20 give \(fx_\beta=y_\alpha\).  The
   deduction theorem and rule E (on \(g\), bound in (b)) give
   (b) \(\derives\mathbf Ey_\alpha\supset fx_\beta=y_\alpha\) after Rule~III, and
   VI\('\) on \(y_\alpha\), not free in (b), gives the claim.
 \item The instance of 9\(^\alpha\) with \(\mathbf E\) for \(f_{o\alpha}\) and
   \(fx_\beta\) for \(x_\alpha\), reduced, and Rule~V twice with items 1 and 2
   give \(\mathbf E(\iota_{\alpha(o\alpha)}\mathbf E)\), that is
   \(\mathbf E(\mathbf Gx_\beta)\) up to conversion.  Item 2 at
   \(\mathbf Gx_\beta\) gives \(fx_\beta=\mathbf Gx_\beta\).
 \item VI\('\) on \(x_\beta\), which is free in neither assumption, gives
   \((x_\beta)[fx_\beta=\mathbf Gx_\beta]\).  The instance of 10\(^{\alpha\beta}\)
   with \(\mathbf G\) for \(g_{\alpha\beta}\), reduced, gives \(f=\mathbf G\).
   Then 17 at type \(\alpha\beta\) with \(h\) for its predicate and (a) give
   \(h\mathbf G\), and Rule~III gives
   \(h(\iota_{\alpha\beta(o(\alpha\beta))}h)\).
\end{enumerate}
The deduction theorem twice gives 9\(^{\alpha\beta}\).  Besides 9\(^\alpha\) and
10\(^{\alpha\beta}\), only Axioms 1--6 (through 12--20 and the deduction
theorem) were used; Axioms 7 and 8 were not needed.

\emph{Clause 2.}  Fix \(\alpha\) and a choice family \(c\) as in
\cref{def:p02-w02-p02-040-church-standard} with \(D=\N\).  Pick
\(a\in\mathcal D_\alpha\) and let \(c'\) agree with \(c\) except that
\(c'_\alpha\) is the constant function with value \(a\).  In the proof of
\cref{thm:p02-w02-p02-040-church-soundness}, the clauses for the rules, for
equality, for Axioms 1--8 and 10, and for 9\(^\beta\) with \(\beta\neq\alpha\)
use no property of \(c_\alpha\); so all of these are valid in the model
\(\mathcal S(\N,c')\) and the rules preserve validity there.  Every
\(\mathcal D_\gamma\) has at least two members (\(\mathcal D_o\) has two,
\(\mathcal D_\iota=\N\), and \(\mathcal D_{\gamma\delta}\) has at least as many
members as \(\mathcal D_\gamma\)), so choose \(b\neq a\) in \(\mathcal D_\alpha\)
and let \(s(f_{o\alpha})\) be the characteristic function of \(\{b\}\) and
\(s(x_\alpha)=b\).  Then \(f_{o\alpha}x_\alpha\) and
\((y_\alpha)[f_{o\alpha}y_\alpha\supset x_\alpha=y_\alpha]\) are true under \(s\),
by the equality clause, while
\(f_{o\alpha}(\iota_{\alpha(o\alpha)}f_{o\alpha})\) has the value of the
characteristic function at \(a\), which is \(\mathrm F\).  So 9\(^\alpha\) is not
valid in \(\mathcal S(\N,c')\), and it is not derivable from the other axioms
listed.
\end{proof}
\end{annexitem}

\begin{annexitem}{annex:p02-w02-p02-040-church-bridge}{thm:p02-w02-p02-040-church-bridge}{The bridge preserves types and denotations}
\begin{proof}
The host type grammar is free, so \(\beta\to\alpha\mapsto(\alpha^\bullet\beta^\bullet)\),
with \(\iota^\bullet=\iota\) and \(o^\bullet=o\), is well defined by recursion,
and induction shows that it is inverse to \(\gamma\mapsto\gamma^\circ\).
Typing of \(\mathbf M^\circ\) is proved by induction on \(\mathbf M\): a
Church variable or constant of type \(\gamma\) goes to a host symbol of
type \(\gamma^\circ\) (for instance \(\Pi_{o(o\alpha)}^\circ=\forall_{\alpha^\circ}\)
has type \((\alpha^\circ\to o)\to o=(o(o\alpha))^\circ\)); if
\(\mathbf F:\alpha\beta\) and \(\mathbf A:\beta\) then
\(\mathbf F^\circ:\beta^\circ\to\alpha^\circ\) and
\(\mathbf A^\circ:\beta^\circ\), so the application has type \(\alpha^\circ\);
and \((\lambda\mathbf x_\beta\mathbf M_\alpha)^\circ\) has type
\(\beta^\circ\to\alpha^\circ=(\alpha\beta)^\circ\).  The translation acts
symbol by symbol on variables and \(\lambda\), so free and bound
occurrences correspond.

By induction on types, the host standard domain \(D_{\gamma^\circ}\) is
\(\mathcal D_\gamma\) once \(1,0\) are identified with \(\mathrm T,\mathrm F\):
both are \(D\) at \(\iota\), the truth values at \(o\), and the set of all
functions at function types.  The host constants \(\neg,\lor,\forall_{\alpha^\circ}\)
have the denotations fixed for \(N_{oo},A_{ooo},\Pi_{o(o\alpha)}\) in
\cref{def:p02-w02-p02-040-church-standard}, and \(\varepsilon_\alpha\) denotes
\(c_\alpha\) by hypothesis.  The variable, application and abstraction
clauses of the two denotation recursions coincide, so induction on
\(\mathbf M\) gives \(\sem{\mathbf M^\circ}^{\mathcal M}_s=\sem{\mathbf M}_s\).

For the defined symbols: \([\mathbf A\supset\mathbf B]^\circ\) is
\(\lor(\neg\mathbf A^\circ)\mathbf B^\circ\), whose value is that of the material
implication \(\to\mathbf A^\circ\mathbf B^\circ\); the conjunction is
\(\neg(\lor(\neg\mathbf A^\circ)(\neg\mathbf B^\circ))\), whose value is that of
\(\land\mathbf A^\circ\mathbf B^\circ\); and
\([(\exists\mathbf x_\alpha)\mathbf A]^\circ\) is
\(\neg\forall_{\alpha^\circ}(\lambda\mathbf x.\neg\mathbf A^\circ)\), which has
the value of \(\exists_{\alpha^\circ}(\lambda\mathbf x.\mathbf A^\circ)\) as
defined in \cref{def:p02-stt-boolean-constants}, by the beta identity of
\cref{sec:simple-type-theory-02}.  Finally
\([\mathbf A_\alpha=\mathbf B_\alpha]^\circ\) has the Church value of
\([\mathbf A_\alpha=\mathbf B_\alpha]\), which by the equality clause of
\cref{thm:p02-w02-p02-040-church-soundness} is \(\mathrm T\) iff the two
denotations are equal, that is, the value of
\(\mathbf A^\circ=_{\alpha^\circ}\mathbf B^\circ\).
\end{proof}
\end{annexitem}

\begin{annexitem}{annex:p02-w02-p02-040-church-general-soundness}{thm:p02-w02-p02-040-church-general-soundness}{Soundness for general models}
\begin{proof}
Fix a general model \(\mathcal M\).

\emph{Denotation lemmas.}  The coincidence, substitution, renaming and
replacement lemmas at the start of
\cref{annex:p02-w02-p02-040-church-soundness} hold in \(\mathcal M\).  Each of
them is proved there by induction on wffs, and each inductive step uses
only the three denotation clauses and the existence of
\(\sem{\mathbf M}_s\) for every wff and assignment; it uses no property of the
domains.  In \(\mathcal M\) the clauses are the same and existence is (G4),
so those inductions prove the four lemmas for \(\mathcal M\).

\emph{Equality.}  With abbreviations eliminated, \([\mathbf A_\alpha=\mathbf B_\alpha]\)
is \(Q_{o\alpha\alpha}\mathbf A_\alpha\mathbf B_\alpha\) with the closed wff
\(Q_{o\alpha\alpha}=\lambda x_\alpha\lambda y_\alpha\Pi_{o(o\alpha)}(\lambda f_{o\alpha}A_{ooo}(N_{oo}(f_{o\alpha}x_\alpha))(f_{o\alpha}y_\alpha))\).
By the denotation clauses and coincidence, with
\(e=\sem{\mathbf A_\alpha}_s\) and \(e'=\sem{\mathbf B_\alpha}_s\),
\(\sem{Q_{o\alpha\alpha}\mathbf A_\alpha\mathbf B_\alpha}_s=\pi_\alpha(h\mapsto a(n(h(e)))(h(e')))\),
which by (G3) lies in \(\mathcal T\) iff every \(h\in\mathcal D_{o\alpha}\) with
\(h(e)\in\mathcal T\) has \(h(e')\in\mathcal T\).  This holds when \(e=e'\).  When
\(e\neq e'\), (G5) supplies an \(h\) with \(h(e)\in\mathcal T\) and
\(h(e')\notin\mathcal T\), so it fails.

\emph{Axioms.}  By (G3) for \(n\) and \(a\), the value of a wff built by
\(\sim\) and \(\vee\) from wffs of type \(o\) lies in \(\mathcal T\) exactly when
the truth table, applied to the memberships of the components in
\(\mathcal T\), gives truth; Axioms 1--4 are tautologies, so they are valid.
Axiom 5\(^\alpha\): if \(\pi_\alpha(s(f_{o\alpha}))\in\mathcal T\) then
\(s(f_{o\alpha})(e)\in\mathcal T\) for every \(e\), in particular for
\(e=s(x_\alpha)\).  Axiom 6\(^\alpha\): the antecedent denotes
\(\pi_\alpha(e\mapsto a(s(p))(s(f_{o\alpha})(e)))\), which lies in \(\mathcal T\)
iff for every \(e\), \(s(p)\in\mathcal T\) or \(s(f_{o\alpha})(e)\in\mathcal T\);
then either \(s(p)\in\mathcal T\) or \(\pi_\alpha(s(f_{o\alpha}))\in\mathcal T\).
Axioms 9\(^\alpha\) and 11\(^\alpha\): if \(s(f_{o\alpha})(s(x_\alpha))\in\mathcal T\),
then \(s(f_{o\alpha})(c_\alpha(s(f_{o\alpha})))\in\mathcal T\) by (G3); the
consequent of 9\(^\alpha\) and that of 11\(^\alpha\) are this statement.
Axiom 10\(^{\alpha\beta}\): by the equality clause the antecedent says that
\(s(f_{\alpha\beta})(e)=s(g_{\alpha\beta})(e)\) for every \(e\in\mathcal D_\beta\);
both are functions on \(\mathcal D_\beta\) by (G1), so they are equal, and
the consequent holds by the equality clause.

\emph{Proofs on assumptions.}  Suppose Axioms 7 and 8 are valid in
\(\mathcal M\), let \(\mathbf C^1_o,\ldots,\mathbf C^m_o\) be a proof of
\(\mathbf B_o\) on \(\mathbf A^1_o,\ldots,\mathbf A^n_o\), and let \(s\) satisfy
every \(\mathbf A^j_o\).  We show by induction on \(i\) that
\(\mathcal M,s'\entails\mathbf C^i_o\) for every assignment \(s'\) that satisfies
every \(\mathbf A^j_o\); the case \(i=m\), \(s'=s\) is the claim.  Assumptions
are satisfied by hypothesis.  Every axiom is valid, and a variant of an
axiom is obtained from the axiom by Rules I and IV\('\)
(\cref{def:p02-w02-p02-040-church-axioms}), which preserve validity by the
cases for Rules I--IV below taken with \(n=0\).  Rules I--III preserve
denotation under every assignment, by renaming, substitution and replacement,
as in \cref{annex:p02-w02-p02-040-church-soundness}.  Rule~IV: let
\(\mathbf C^i_o=\mathbf F\mathbf A\) come from \(\mathbf F\mathbf x\), where
\(\mathbf x\) is free neither in \(\mathbf F\) nor, by the restriction, in any
assumption.  Given \(s'\), the assignment
\(s''=s'[\mathbf x\mapsto\sem{\mathbf A}_{s'}]\) still satisfies the assumptions
(coincidence), so \(\sem{\mathbf F\mathbf x}_{s''}\in\mathcal T\) by the
induction hypothesis, and
\(\sem{\mathbf F\mathbf x}_{s''}=\sem{\mathbf F}_{s'}(\sem{\mathbf A}_{s'})=\sem{\mathbf F\mathbf A}_{s'}\).
Rule~V: by (G3), if \(\sem{\mathbf C\supset\mathbf D}_{s'}\) and
\(\sem{\mathbf C}_{s'}\) lie in \(\mathcal T\), so does \(\sem{\mathbf D}_{s'}\).
Rule~VI: let \(\Pi_{o(o\alpha)}\mathbf F\) come from \(\mathbf F\mathbf x\), with
\(\mathbf x\) free neither in \(\mathbf F\) nor in any assumption.  For every
\(e\in\mathcal D_\alpha\), \(s'[\mathbf x\mapsto e]\) satisfies the assumptions,
so \(\sem{\mathbf F}_{s'}(e)=\sem{\mathbf F\mathbf x}_{s'[\mathbf x\mapsto e]}\in\mathcal T\);
hence \(\pi_\alpha(\sem{\mathbf F}_{s'})\in\mathcal T\).  The last sentence of
the theorem is the case \(n=0\).
\end{proof}
\end{annexitem}

\begin{annexitem}{annex:p02-w02-p02-040-church-encoding}{thm:p02-w02-p02-040-church-encoding}{Encoding and decoding}
\begin{proof}
\emph{Clause 1.}  By (G4) and coincidence,
\(\ell_{\mathbf L}^{\mathcal M^\dagger}(e_1,\ldots,e_k)\) lies in
\(\mathcal D_{\alpha\beta}\) and equals \(\sem{\mathbf L}_s\) for every \(s\) with
\(s(\mathbf y^i)=e_i\).  Induction on \(\mathbf M\) shows that the value of
\(\mathbf M^\dagger\) under \(s\) is \(\sem{\mathbf M}_s\): variables, constants
and applications match clause by clause, and for an abstract \(\mathbf L\)
the value of \(\ell_{\mathbf L}(\mathbf y^1,\ldots,\mathbf y^k)\) under \(s\) is
\(\sem{\mathbf L}_s\) by the remark just made.  Now check \(\mathsf{Enc}\).
\((\mathrm E_N)\), \((\mathrm E_A)\), \((\mathrm E_\Pi^\alpha)\) and
\((\mathrm E_\iota^\alpha)\) are the four conditions of (G3).
\((\mathrm E_\beta^{\mathbf L})\): under \(s\) the left side is
\(\sem{\mathbf L}_s(s(\mathbf x_\beta))=\sem{\mathbf M_\alpha}_{s[\mathbf x_\beta\mapsto s(\mathbf x_\beta)]}=\sem{\mathbf M_\alpha}_s\),
which is the value of the right side.  In \(\mathcal M^\dagger\),
\(\mathsf{Lb}_\gamma(u,v)\) says that every \(h\in\mathcal D_{o\gamma}\) with
\(h(u)\in\mathcal T\) has \(h(v)\in\mathcal T\), which by (G5) holds iff \(u=v\).
So \((\mathrm E_{\mathrm{ext}}^{\alpha\beta})\) says that two members of
\(\mathcal D_{\alpha\beta}\) that agree at every argument are equal, which is
true of functions.  \((\mathrm E_7)\) and \((\mathrm E_8)\) hold because Axioms 7
and 8 are valid and values of encodings are denotations.  The final
equivalence follows from the value identity.

\emph{Clause 2, preliminaries.}  Let \(\mathfrak A\models\mathsf{Enc}\).  Write
\(U_\gamma\) for the domain of \(\mathsf s_\gamma\), \(u\cdot w\) for
\(\mathsf{ap}^{\mathfrak A}(u,w)\), \(\mathrm{val}(\mathbf M,t)\) for the value
of \(\mathbf M^\dagger\) under \(t\), and \(u\sim_\gamma v\) for
\(\mathfrak A\models\mathsf{Lb}_\gamma(u,v)\).  Instantiating
\((\mathrm E_\beta^{\mathbf L})\) gives, for every abstract
\(\mathbf L=\lambda\mathbf x\mathbf M\), every \(t\) and every \(w\),
\begin{equation}
 \mathrm{val}(\mathbf L,t)\cdot w=\mathrm{val}(\mathbf M,t[\mathbf x\mapsto w]).
 \tag{\(\beta\)}
\end{equation}
\emph{Leibniz encoding.}  For all \(u,v\in U_\alpha\),
\(\mathsf{True}(\mathrm{val}(Q_{o\alpha\alpha},t)\cdot u\cdot v)\) iff
\(u\sim_\alpha v\).  Indeed \(Q_{o\alpha\alpha}\) is the abstract
\(\lambda x_\alpha\mathbf L_2\) with
\(\mathbf L_2=\lambda y_\alpha\Pi_{o(o\alpha)}\mathbf L_3\) and
\(\mathbf L_3=\lambda f_{o\alpha}A_{ooo}(N_{oo}(f_{o\alpha}x_\alpha))(f_{o\alpha}y_\alpha)\).
Three uses of (\(\beta\)) give
\(\mathrm{val}(Q,t)\cdot u\cdot v=\widehat\Pi^{\mathfrak A}_\alpha\cdot r\) with
\(r\cdot h=\widehat A\cdot(\widehat N\cdot(h\cdot u))\cdot(h\cdot v)\) for every
\(h\), where \(r=\mathrm{val}(\mathbf L_3,t[x_\alpha\mapsto u,y_\alpha\mapsto v])\);
by \((\mathrm E_\Pi^\alpha)\), \((\mathrm E_A)\) and \((\mathrm E_N)\) the
proposition is true iff for every \(h\), \(\mathsf{True}(h\cdot u)\) fails or
\(\mathsf{True}(h\cdot v)\) holds.

\emph{Clause 2, Step 1: \(\sim\) is a congruence.}  Reflexivity and
transitivity are immediate from the form of \(\mathsf{Lb}\).  For symmetry, let
\(u\sim v\) and let \(h_0=\mathrm{val}(\lambda z_\gamma(Q_{o\gamma\gamma}z_\gamma w_\gamma),t[w_\gamma\mapsto u])\).
By (\(\beta\)) and the Leibniz encoding, \(\mathsf{True}(h_0\cdot z)\) iff
\(z\sim u\).  Since \(u\sim u\), \(\mathsf{True}(h_0\cdot u)\), so
\(\mathsf{True}(h_0\cdot v)\), that is, \(v\sim u\).  Truth respects \(\sim_o\):
with \(h=\mathrm{val}(\lambda p_o\,p_o,t)\), (\(\beta\)) gives \(h\cdot u=u\), so
\(u\sim_o v\) and \(\mathsf{True}(u)\) give \(\mathsf{True}(v)\).  Application
respects \(\sim\) in each argument.  If \(w\sim_\beta w'\) and
\(u\in U_{\alpha\beta}\), then for any \(h\in U_{o\alpha}\) let
\(h'=\mathrm{val}(\lambda z_\beta(g_{o\alpha}(f_{\alpha\beta}z_\beta)),t[g\mapsto h,f\mapsto u])\);
by (\(\beta\)), \(h'\cdot z=h\cdot(u\cdot z)\), and \(w\sim w'\) applied to \(h'\)
gives \(u\cdot w\sim u\cdot w'\).  If \(u\sim_{\alpha\beta}u'\) and
\(w\in U_\beta\), let
\(h'=\mathrm{val}(\lambda f_{\alpha\beta}(g_{o\alpha}(f_{\alpha\beta}z_\beta)),t[g\mapsto h,z\mapsto w])\);
then \(h'\cdot f=h\cdot(f\cdot w)\), and \(u\sim u'\) gives
\(u\cdot w\sim u'\cdot w\).

\emph{Step 2: the frame.}  Define \(\mathcal D_\gamma\) and
\(d_\gamma:U_\gamma\to\mathcal D_\gamma\) by recursion on \(\gamma\).  For
\(\gamma=\iota,o\) let \(\mathcal D_\gamma=U_\gamma/{\sim_\gamma}\) and
\(d_\gamma(u)=[u]\); let \(\mathcal T=\{d_o(u):\mathsf{True}(u)\}\).  For
\(\gamma=\alpha\beta\) let \(\varphi_u:\mathcal D_\beta\to\mathcal D_\alpha\) be
\(\varphi_u(d_\beta(w))=d_\alpha(u\cdot w)\), let \(d_{\alpha\beta}(u)=\varphi_u\)
and \(\mathcal D_{\alpha\beta}=\{\varphi_u:u\in U_{\alpha\beta}\}\).  We prove by
induction on \(\gamma\) that \(d_\gamma\) is a surjection onto
\(\mathcal D_\gamma\), that each \(\varphi_u\) is a well-defined function, and
that \(d_\gamma(u)=d_\gamma(v)\) iff \(u\sim_\gamma v\).  At base types this is
the definition.  At \(\alpha\beta\): \(d_\beta\) is onto, so \(\varphi_u\) is
defined on all of \(\mathcal D_\beta\); it is well defined because
\(d_\beta(w)=d_\beta(w')\) gives \(w\sim w'\) (induction hypothesis), hence
\(u\cdot w\sim u\cdot w'\) (Step~1), hence
\(d_\alpha(u\cdot w)=d_\alpha(u\cdot w')\).  If \(u\sim u'\), Step~1 gives
\(u\cdot w\sim u'\cdot w\) for all \(w\), so \(\varphi_u=\varphi_{u'}\).
Conversely, if \(\varphi_u=\varphi_{u'}\), then \(u\cdot w\sim u'\cdot w\) for
every \(w\) by the induction hypothesis, and
\((\mathrm E_{\mathrm{ext}}^{\alpha\beta})\) gives \(u\sim u'\).  Thus (G1) holds,
the domains being nonempty because sorts are; (G2) holds, and
\(d_o(u)\in\mathcal T\) iff \(\mathsf{True}(u)\) because truth respects
\(\sim_o\).

\emph{Step 3: constants, (G3) and (G5).}  Let \(n=d(\widehat N^{\mathfrak A})\),
\(a=d(\widehat A^{\mathfrak A})\), \(\pi_\alpha=d(\widehat\Pi_\alpha^{\mathfrak A})\),
\(c_\alpha=d(\widehat\iota_\alpha^{\mathfrak A})\).  For \(u\in U_o\),
\(n(d(u))=d(\widehat N\cdot u)\), which lies in \(\mathcal T\) iff
\(\mathsf{True}(\widehat N\cdot u)\) iff not \(\mathsf{True}(u)\)
\((\mathrm E_N)\) iff \(d(u)\notin\mathcal T\).  The clause for \(a\) follows in
the same way from \((\mathrm E_A)\).  For \(g=\varphi_u\in\mathcal D_{o\alpha}\),
\(\pi_\alpha(g)=d(\widehat\Pi_\alpha\cdot u)\in\mathcal T\) iff
\(\mathsf{True}(u\cdot w)\) for every \(w\in U_\alpha\) \((\mathrm E_\Pi^\alpha)\)
iff \(g(e)\in\mathcal T\) for every \(e\in\mathcal D_\alpha\), since \(d_\alpha\) is
onto.  If \(g(d(w))\in\mathcal T\), then \(\mathsf{True}(u\cdot w)\), so
\(\mathsf{True}(u\cdot(\widehat\iota_\alpha\cdot u))\) by
\((\mathrm E_\iota^\alpha)\), that is, \(g(c_\alpha(g))\in\mathcal T\).  For (G5),
if \(d(u)\neq d(v)\) then \(u\not\sim v\), so some \(h\) has
\(\mathsf{True}(h\cdot u)\) and not \(\mathsf{True}(h\cdot v)\); then
\(\varphi_h\in\mathcal D_{o\alpha}\) separates \(d(u)\) from \(d(v)\).

\emph{Step 4: closure and denotation.}  We prove by induction on
\(\mathbf M\): for every assignment \(s\) into the frame and every
\(\mathfrak A\)-assignment \(t\) with \(d(t(\mathbf y))=s(\mathbf y)\) for each
free variable \(\mathbf y\) of \(\mathbf M\), the denotation
\(\sem{\mathbf M}_s\) exists and equals \(d(\mathrm{val}(\mathbf M,t))\).
Variables and constants are immediate.  For \(\mathbf F\mathbf A\),
\(\sem{\mathbf F\mathbf A}_s=\varphi_{\mathrm{val}(\mathbf F,t)}(d(\mathrm{val}(\mathbf A,t)))=d(\mathrm{val}(\mathbf F,t)\cdot\mathrm{val}(\mathbf A,t))=d(\mathrm{val}(\mathbf F\mathbf A,t))\).
For \(\mathbf L=\lambda\mathbf x_\beta\mathbf M\) and \(e\in\mathcal D_\beta\), pick
one \(w\) with \(d(w)=e\); then \(t[\mathbf x\mapsto w]\) and
\(s[\mathbf x\mapsto e]\) satisfy the hypothesis for \(\mathbf M\), so by the
induction hypothesis and (\(\beta\))
\(\sem{\mathbf M}_{s[\mathbf x\mapsto e]}=d(\mathrm{val}(\mathbf M,t[\mathbf x\mapsto w]))=d(\mathrm{val}(\mathbf L,t)\cdot w)=\varphi_{\mathrm{val}(\mathbf L,t)}(e)\).
So the function \(e\mapsto\sem{\mathbf M}_{s[\mathbf x\mapsto e]}\) is
\(\varphi_{\mathrm{val}(\mathbf L,t)}\in\mathcal D_{\alpha\beta}\), and
\(\sem{\mathbf L}_s=d(\mathrm{val}(\mathbf L,t))\).  Every \(s\) has such a \(t\):
modify \(t_0\) at the finitely many free variables of \(\mathbf M\), choosing
one preimage under \(d\) for each.  Hence (G4) holds, \(\mathfrak A^\natural\) is
a general model, and \(\sem{\mathbf M}_{d\circ t}=d(\mathrm{val}(\mathbf M,t))\)
for every \(t\).  The truth equivalence follows from
\(d_o(u)\in\mathcal T\Leftrightarrow\mathsf{True}(u)\).  Axioms 7 and 8 are valid:
for any \(s\), take \(t\) as above for the wff concerned; \((\mathrm E_7)\) and
\((\mathrm E_8)\) give \(\mathsf{True}(\mathrm{val}(\cdot,t))\), so the denotation
under \(s\) lies in \(\mathcal T\).

\emph{Clause 3.}  In \(\mathcal M^\dagger\), \(\sim_\gamma\) is identity (clause
1), so each \(d_\gamma\) is injective, hence bijective.  For \(f\in\mathcal D_{\alpha\beta}\)
and \(e\in\mathcal D_\beta\),
\(d_{\alpha\beta}(f)(d_\beta(e))=d_\alpha(f\cdot e)=d_\alpha(f(e))\), and the
constants and \(\mathcal T\) correspond by the definitions in Step~3.
\end{proof}
\end{annexitem}

\begin{annexitem}{annex:p02-w02-p02-040-church-general-completeness}{thm:p02-w02-p02-040-church-general-completeness}{Strong completeness for general models}
\begin{proof}
\emph{1 implies 2.}  If \(\Gamma_0\derives\mathbf B_o\) with \(\Gamma_0\) a finite
list from \(\Gamma\), then \(\Gamma_0\entails_{\mathrm G}\mathbf B_o\) by
\cref{thm:p02-w02-p02-040-church-general-soundness}, and so
\(\Gamma\entails_{\mathrm G}\mathbf B_o\).

\emph{2 implies 3.}  Let \(\mathfrak A,t\) satisfy \(\mathsf{Enc}\) and every
\(\mathsf{True}(\mathbf A_o^\dagger)\), \(\mathbf A_o\in\Gamma\).  By clause~2 of
\cref{thm:p02-w02-p02-040-church-encoding}, with \(t_0=t\), the general model
\(\mathfrak A^\natural\), in which Axioms 7 and 8 are valid, satisfies every
member of \(\Gamma\) under \(d\circ t\).  By 2 it satisfies \(\mathbf B_o\), so
\(\mathfrak A,t\entails\mathsf{True}(\mathbf B_o^\dagger)\).

\emph{3 implies 1: tools.}  Write \(\Delta_0\derives\mathbf C\) for derivability
on a finite list \(\Delta_0\).  The following facts are used.
\begin{enumerate}
 \item[(T1)] If \(\Delta_0\derives\mathbf C\) and every member of \(\Delta_0\)
   occurs in the finite list \(\Delta_1\), then \(\Delta_1\derives\mathbf C\).
   Apply \cref{thm:p02-w02-p02-040-church-deduction} repeatedly to obtain a
   theorem \(\mathbf A^1\supset.\,\cdots\supset.\,\mathbf A^n\supset\mathbf C\),
   use it on \(\Delta_1\) by Step~1 of
   \cref{annex:p02-w02-p02-040-church-deduction}, and detach the
   assumptions by Rule~V.
 \item[(T2)] Two proofs on the same list can be concatenated; so
   \(\Delta_0\derives\mathbf C\) and \(\Delta_0\derives\mathbf C\supset\mathbf D\)
   give \(\Delta_0\derives\mathbf D\), and tautology instances, theorems by
   \cref{thm:p02-w02-p02-040-church-prop-complete}, are available on every
   list by Step~1 of \cref{annex:p02-w02-p02-040-church-deduction}.
 \item[(T3)] (Substitution in theorems.)  If \(\derives\mathbf M_o\) and
   \(\theta\) assigns to finitely many variables \(\mathbf y\) wffs
   \(\theta(\mathbf y)\) of the same types, let \(\mathbf M\theta\) be the result of
   renaming the bound variables of \(\mathbf M\) (Rule~I) to variables occurring
   in no \(\theta(\mathbf y)\) and then replacing each free \(\mathbf y\) by
   \(\theta(\mathbf y)\) simultaneously.  Then \(\derives\mathbf M_o\theta\):
   substitute fresh distinct variables \(\mathbf z\) for the \(\mathbf y\) by
   IV\('\), then \(\theta(\mathbf y)\) for each \(\mathbf z\) by IV\('\), one at a
   time; since no \(\mathbf z\) occurs in any \(\theta(\mathbf y)\), the result is
   the simultaneous substitution.
 \item[(T4)] (Conversion and substitution.)  Write
   \(\mathbf M\mathrel{\mathrm{cnv}}\mathbf M'\) when \(\mathbf M'\) comes from \(\mathbf M\) by finitely
   many applications of Rules I--III.  This is an equivalence relation
   (Rule~III is the converse of Rule~II, and a renaming by Rule~I is undone by
   Rule~I), and it is compatible with forming applications, since a part of
   \(\mathbf F\) or of \(\mathbf A\) is a part of \(\mathbf F\mathbf A\).  For
   \(\theta\) as in (T3), \((\mathbf F\mathbf A)\theta\mathrel{\mathrm{cnv}}(\mathbf F\theta)(\mathbf A\theta)\)
   and \(\mathbf y\theta=\theta(\mathbf y)\); if \(\theta(\mathbf y)\mathrel{\mathrm{cnv}}\theta'(\mathbf y)\)
   for every \(\mathbf y\), then \(\mathbf M\theta\mathrel{\mathrm{cnv}}\mathbf M\theta'\), by converting
   inside the inserted parts; and for \(\mathbf L=\lambda\mathbf x\mathbf M\) and any
   \(\mathbf C\), \((\mathbf L\theta)\mathbf C\mathrel{\mathrm{cnv}}\mathbf M(\theta[\mathbf x\mapsto\mathbf C])\),
   because \(\mathbf L\theta\) is \(\lambda\mathbf x'\,\mathbf M(\theta[\mathbf x\mapsto\mathbf x'])\)
   with \(\mathbf x'\) occurring in no \(\theta(\mathbf y)\); after Rule~I renames
   the bound variables of this body away from \(\mathbf x'\) and from the free
   variables of \(\mathbf C\), one application of Rule~II replaces \(\mathbf x'\)
   by \(\mathbf C\), and the result differs from
   \(\mathbf M(\theta[\mathbf x\mapsto\mathbf C])\) only in the names of bound
   variables.
\end{enumerate}
A finite list \(\Delta_0\) is \emph{inconsistent} if \(\Delta_0\derives\mathbf C\)
and \(\Delta_0\derives\,\sim\mathbf C\) for some \(\mathbf C\); a set is
\emph{consistent} if no finite list of its members is inconsistent.  By (T2)
and the tautology \(q\supset.\,\sim q\supset r\), an inconsistent list derives
every wff of type \(o\).
\begin{enumerate}
 \item[(C1)] If \(\Gamma\cup\{\sim\mathbf B\}\) is inconsistent, then some finite
   \(\Gamma_0\subseteq\Gamma\) has \(\Gamma_0\derives\mathbf B\).  By (T1) take an
   inconsistent list \(\Gamma_0,\sim\mathbf B\); it derives \(\mathbf B\); the
   deduction theorem gives \(\Gamma_0\derives\,\sim\mathbf B\supset\mathbf B\), and
   the tautology \((\sim q\supset q)\supset q\) with (T2) gives
   \(\Gamma_0\derives\mathbf B\).
 \item[(C2)] If \(\Delta\) is consistent, then \(\Delta\cup\{\mathbf C\}\) or
   \(\Delta\cup\{\sim\mathbf C\}\) is consistent.  Otherwise (T1) gives one list
   \(\Delta_0\) of members of \(\Delta\) with \(\Delta_0,\mathbf C\) and
   \(\Delta_0,\sim\mathbf C\) inconsistent.  Then
   \(\Delta_0\derives\mathbf C\supset\,\sim\mathbf C\) and
   \(\Delta_0\derives\,\sim\mathbf C\supset\mathbf C\) by the deduction theorem,
   and the tautologies \((q\supset\,\sim q)\supset\,\sim q\) and
   \((\sim q\supset q)\supset q\) make \(\Delta_0\) inconsistent.
 \item[(C3)] (Witnesses.)  If \(\Delta\) is consistent, \(\mathbf F_{o\alpha}\) is a
   wff and \(\mathbf y_\alpha\) is free neither in \(\mathbf F_{o\alpha}\) nor in any
   member of \(\Delta\), then \(\Delta\cup\{\mathbf F\mathbf y\supset\Pi_{o(o\alpha)}\mathbf F\}\)
   is consistent.  Otherwise some list \(\Delta_0,\mathbf W\) with
   \(\mathbf W=\mathbf F\mathbf y\supset\Pi\mathbf F\) is inconsistent; as in (C2),
   \(\Delta_0\derives\,\sim\mathbf W\), and the tautologies
   \(\sim(q\supset r)\supset q\) and \(\sim(q\supset r)\supset\,\sim r\) give
   \(\Delta_0\derives\mathbf F\mathbf y\) and \(\Delta_0\derives\,\sim\Pi\mathbf F\).
   Appending Rule~VI to the first proof is allowed, because \(\mathbf y\) is free
   neither in \(\mathbf F\) nor in any assumption; so \(\Delta_0\derives\Pi\mathbf F\)
   and \(\Delta_0\) is inconsistent.
\end{enumerate}

\emph{3 implies 1: a special case.}  Suppose first that infinitely many
variables of each type are free in no member of \(\Gamma\cup\{\mathbf B\}\),
and that 1 fails.  By (C1), \(\Delta_0\defeq\Gamma\cup\{\sim\mathbf B\}\) is
consistent.  The wffs of type \(o\) are finite strings over a countable
alphabet; fix an enumeration \(\mathbf C_0,\mathbf C_1,\ldots\) of them.  Given
consistent \(\Delta_k\), in which infinitely many variables of each type are
not free, let \(\Delta_k'=\Delta_k\cup\{\mathbf F\mathbf y\supset\mathbf C_k\}\) if
\(\mathbf C_k\) has the form \(\Pi_{o(o\alpha)}\mathbf F_{o\alpha}\), where
\(\mathbf y\) is the first variable of type \(\alpha\) in the enumeration free in
no member of \(\Delta_k\cup\{\mathbf C_k\}\), and \(\Delta_k'=\Delta_k\) otherwise;
then let \(\Delta_{k+1}=\Delta_k'\cup\{\mathbf C_k\}\) if this is consistent, and
\(\Delta_k'\cup\{\sim\mathbf C_k\}\) otherwise.  By (C3) and (C2) every
\(\Delta_k\) is consistent, and each step adds at most two wffs, so the supply
of unused variables persists.  Let \(\Delta^*=\bigcup_k\Delta_k\); it is
consistent, since a finite list lies in some \(\Delta_k\).  It has the
following properties, for all wffs.
\begin{enumerate}
 \item[(M1)] Exactly one of \(\mathbf C\), \(\sim\mathbf C\) lies in \(\Delta^*\).
 \item[(M2)] If a finite list of members derives \(\mathbf C\), then
   \(\mathbf C\in\Delta^*\); otherwise \(\sim\mathbf C\in\Delta^*\) and the list with
   \(\sim\mathbf C\) added is inconsistent.  So \(\Delta^*\) contains every
   theorem (T1), is closed under Rules I--III and V, and if
   \(\mathbf C\mathrel{\mathrm{cnv}}\mathbf C'\) then \(\mathbf C\in\Delta^*\) iff \(\mathbf C'\in\Delta^*\).
 \item[(M3)] \(\mathbf C\vee\mathbf D\in\Delta^*\) iff \(\mathbf C\in\Delta^*\) or
   \(\mathbf D\in\Delta^*\), by (M1), (M2) and the tautologies
   \(q\supset q\vee r\), \(r\supset q\vee r\) and
   \(q\vee r\supset.\,\sim q\supset r\).
 \item[(M4)] \(\Pi_{o(o\alpha)}\mathbf F\in\Delta^*\) iff
   \(\mathbf F\mathbf A\in\Delta^*\) for every wff \(\mathbf A_\alpha\).  From left to
   right use the theorem \(\Pi\mathbf F\supset\mathbf F\mathbf A\), an instance of
   5\(^\alpha\) by (T3).  From right to left, the witness
   \(\mathbf F\mathbf y\supset\Pi\mathbf F\) added when \(\Pi\mathbf F\) was enumerated
   lies in \(\Delta^*\), and \(\mathbf F\mathbf y\in\Delta^*\).
 \item[(M5)] If \(\mathbf F\mathbf A\in\Delta^*\) then
   \(\mathbf F(\iota_{\alpha(o\alpha)}\mathbf F)\in\Delta^*\), by the instance of
   11\(^\alpha\) given by (T3).
 \item[(M6)] \([\mathbf A_\alpha=\mathbf B_\alpha]\in\Delta^*\) iff every
   \(\mathbf H_{o\alpha}\) with \(\mathbf H\mathbf A\in\Delta^*\) has
   \(\mathbf H\mathbf B\in\Delta^*\).  From left to right use the instance of
   17\(^\alpha\) (\cref{thm:p02-w02-p02-040-church-12-21}).  From right to left:
   \(Q\mathbf A\mathbf B\mathrel{\mathrm{cnv}}\Pi(\lambda f[f\mathbf A\supset f\mathbf B])\) by Rule~II,
   with \(f\) free in neither wff; if it is not in \(\Delta^*\), (M4) gives an
   \(\mathbf H\) with \((\lambda f[f\mathbf A\supset f\mathbf B])\mathbf H\notin\Delta^*\),
   hence \([\mathbf H\mathbf A\supset\mathbf H\mathbf B]\notin\Delta^*\) by (M2), hence
   \(\mathbf H\mathbf A\in\Delta^*\) and \(\mathbf H\mathbf B\notin\Delta^*\) by (M1), (M3).
 \item[(M7)] If \([\mathbf F\mathbf A=\mathbf G\mathbf A]\in\Delta^*\) for every
   \(\mathbf A_\beta\), then \([\mathbf F=\mathbf G]\in\Delta^*\).  Take \(\mathbf x_\beta\)
   free in neither.  For each \(\mathbf A\),
   \((\lambda\mathbf x[\mathbf F\mathbf x=\mathbf G\mathbf x])\mathbf A\mathrel{\mathrm{cnv}}[\mathbf F\mathbf A=\mathbf G\mathbf A]\),
   so (M2) and (M4) put \((\mathbf x)[\mathbf F\mathbf x=\mathbf G\mathbf x]\) in
   \(\Delta^*\), and the instance of 10\(^{\alpha\beta}\) given by (T3) yields
   \([\mathbf F=\mathbf G]\).
\end{enumerate}

Now define the \(\Sigma_{\mathrm C}\)-structure \(\mathfrak A\).  The domain of
\(\mathsf s_\gamma\) is the set of \(\mathrel{\mathrm{cnv}}\)-classes \([\mathbf A]\) of wffs of
type \(\gamma\); \([\mathbf F]\cdot[\mathbf A]=[\mathbf F\mathbf A]\) (well defined by
(T4)); the constants are the classes of \(N_{oo}\), \(A_{ooo}\),
\(\Pi_{o(o\alpha)}\), \(\iota_{\alpha(o\alpha)}\); for an abstract \(\mathbf L\) with
free variables \(\mathbf y^1,\ldots,\mathbf y^k\),
\(\ell_{\mathbf L}([\mathbf B^1],\ldots,[\mathbf B^k])=[\mathbf L\theta]\) with
\(\theta(\mathbf y^i)=\mathbf B^i\) (well defined by (T4)); \(\mathsf{True}([\mathbf A])\)
iff \(\mathbf A\in\Delta^*\) (well defined by (M2)); and \(=\) is identity.
Every sort is nonempty.  If \(t\) is an \(\mathfrak A\)-assignment and
\(\theta\) assigns to each free variable \(\mathbf y\) of \(\mathbf M\) a member of
\(t(\mathbf y)\), then the value of \(\mathbf M^\dagger\) under \(t\) is
\([\mathbf M\theta]\): by induction on \(\mathbf M\), using (T4) for applications
and the definition of \(\ell_{\mathbf L}\) for abstracts.

\(\mathfrak A\) is a model of \(\mathsf{Enc}\).  \((\mathrm E_N)\): by (M1)
\(\sim\mathbf A\in\Delta^*\) iff \(\mathbf A\notin\Delta^*\).  \((\mathrm E_A)\) is
(M3), \((\mathrm E_\Pi^\alpha)\) is (M4) and \((\mathrm E_\iota^\alpha)\) is (M5),
because every element of a sort is a class \([\mathbf A]\).
\((\mathrm E_\beta^{\mathbf L})\): for \(\mathbf L=\lambda\mathbf x\mathbf M\) and values
\([\mathbf B^i]\) of the \(\mathbf y^i\) and \([\mathbf C]\) of \(\mathbf x\), the left side
is \([(\mathbf L\theta)\mathbf C]=[\mathbf M(\theta[\mathbf x\mapsto\mathbf C])]\) by (T4),
the value of the right side.  By (M6), \(\mathfrak A\models\mathsf{Lb}([\mathbf A],[\mathbf B])\)
iff \([\mathbf A=\mathbf B]\in\Delta^*\); so \((\mathrm E_{\mathrm{ext}}^{\alpha\beta})\) is
(M7).  \((\mathrm E_7)\) and \((\mathrm E_8)\): the value of \(7^\dagger\) is
\([7]\), and the value of \(8^\dagger\) under any \(t\) is \([8\theta]\); both
are theorems ((T3) for the second), so they lie in \(\Delta^*\) by (M2).
Under \(t_0(\mathbf x)=[\mathbf x]\) the value of \(\mathbf A^\dagger\) is \([\mathbf A]\),
so \(\mathfrak A,t_0\) satisfies \(\mathsf{True}(\mathbf A^\dagger)\) for
\(\mathbf A\in\Gamma\subseteq\Delta^*\) and not \(\mathsf{True}(\mathbf B^\dagger)\),
since \(\sim\mathbf B\in\Delta^*\).  So 3 fails.  No choice was used: the
enumeration, the structure and \(t_0\) are explicitly defined.

\emph{3 implies 1: the general case.}  Let \(\rho\) send the \(j\)-th variable
of each type in the enumeration to the \(2j\)-th one of that type, and
write \(\rho\mathbf M\) for the result of replacing every variable occurrence,
bound ones included, by its image.  Suppose 1 fails for \(\Gamma,\mathbf B\).
Then it fails for \(\rho\Gamma,\rho\mathbf B\): given a proof of \(\rho\mathbf B\) on
\(\rho\mathbf A^1,\ldots,\rho\mathbf A^n\), let \(V\) be the finite set of
variables occurring in it, define \(\sigma\) on \(V\) by
\(\sigma(\rho\mathbf x)=\mathbf x\) for \(\rho\mathbf x\in V\), and send the other
members of \(V\) injectively to variables of the same types occurring
neither in \(\mathbf A^1,\ldots,\mathbf A^n,\mathbf B\) nor among the \(\mathbf x\) with
\(\rho\mathbf x\in V\).  This \(\sigma\) is injective and type-preserving, so,
exactly as in Step~1 of \cref{annex:p02-w02-p02-040-church-deduction}, applying
it throughout the proof gives a proof of \(\mathbf B\) on
\(\mathbf A^1,\ldots,\mathbf A^n\), contrary to assumption.  In
\(\rho\Gamma\cup\{\rho\mathbf B\}\) no variable of odd index occurs, so the
special case gives \(\mathfrak A,t\) satisfying \(\mathsf{Enc}\) and
\(\mathsf{True}((\rho\mathbf A)^\dagger)\) for \(\mathbf A\in\Gamma\), but not
\(\mathsf{True}((\rho\mathbf B)^\dagger)\).  Clause~2 of
\cref{thm:p02-w02-p02-040-church-encoding} with \(t_0=t\) turns this into the
general model \(\mathfrak A^\natural\) and the assignment \(s'=d\circ t\)
satisfying \(\rho\Gamma\) and refuting \(\rho\mathbf B\).  In any general model,
\(\sem{\rho\mathbf M}_{s'}=\sem{\mathbf M}_{s'\circ\rho}\) by induction on
\(\mathbf M\); in the abstraction case
\((s'[\rho\mathbf x\mapsto e])\circ\rho=(s'\circ\rho)[\mathbf x\mapsto e]\) because
\(\rho\) is injective.  So \(s=s'\circ\rho\) satisfies \(\Gamma\) and refutes
\(\mathbf B\) in \(\mathfrak A^\natural\).  Clause~1 of
\cref{thm:p02-w02-p02-040-church-encoding}, with \(s_0=s\), gives the model
\((\mathfrak A^\natural)^\dagger\) of \(\mathsf{Enc}\) in which \(s\) satisfies every
\(\mathsf{True}(\mathbf A^\dagger)\), \(\mathbf A\in\Gamma\), and not
\(\mathsf{True}(\mathbf B^\dagger)\).  So 3 fails.  Each step names its
objects explicitly, so no choice is used.

\emph{Consequences.}  With \(\Gamma=\varnothing\), 1 and 2 say that the
theorems are the wffs valid in every general model in which 7 and 8 are
valid.  If \(\Gamma\entails_{\mathrm G}\mathbf B\), then 1 gives a finite
\(\Gamma_0\subseteq\Gamma\) with \(\Gamma_0\derives\mathbf B\), and soundness gives
\(\Gamma_0\entails_{\mathrm G}\mathbf B\): the relation is compact.
\end{proof}
\end{annexitem}

\begin{annexitem}{annex:p02-w02-p02-040-sol-native}{thm:p02-w02-p02-040-sol-native}{Soundness and completeness of the native relational calculus}
\begin{proof}
Throughout, \(\mathcal H\) is a relational Henkin structure with
\(P^{\mathcal H}\in\mathcal R_n\) for every \(n\)-ary predicate constant, \(s\) a
matching assignment, \(\sem{t}_s\) the value of an object term, and, for an
\(n\)-ary relation variable or predicate constant \(U\),
\(\sem{U}_s=s(U)\) or \(U^{\mathcal H}\) respectively, a member of
\(\mathcal R_n\).  Let \(c(\varphi)\) be the number of occurrences of
connectives and quantifiers in \(\varphi\); renaming a bound variable and the
substitutions \([t/x]\) and \([U/X]\) do not change it.

\emph{Semantic lemmas.}  (L1) If \(s,s'\) agree on the free object and
relation variables of \(\varphi\), then \(\mathcal H,s\entails\varphi\) iff
\(\mathcal H,s'\entails\varphi\).  (L2)
\(\mathcal H,s\entails\varphi[t/x]\) iff
\(\mathcal H,s[x\mapsto\sem t_s]\entails\varphi\).  (L3)
\(\mathcal H,s\entails\varphi[U/X]\) iff
\(\mathcal H,s[X\mapsto\sem U_s]\entails\varphi\).  (L1) is proved by induction
on \(\varphi\), the quantifier steps modifying both assignments at the bound
variable.  (L2) and (L3) are proved together with the renaming facts
\(\mathcal H,s\entails\forall y\,\psi\) iff
\(\mathcal H,s\entails\forall y'\,\psi[y'/y]\) and
\(\mathcal H,s\entails\forall Y\,\psi\) iff
\(\mathcal H,s\entails\forall Y'\,\psi[Y'/Y]\) (and the same for \(\exists\)) for
fresh \(y',Y'\), by induction on \(c(\varphi)\).  For (L3): an atom
\(X(\bar t)\) becomes \(U(\bar t)\), true under \(s\) iff
\(\sem{\bar t}_s\in\sem U_s\), iff \(X(\bar t)\) is true under
\(s[X\mapsto\sem U_s]\); other atoms do not contain \(X\), and (L1) applies.
Connectives commute with the substitution.  Under an object quantifier, or a
relation quantifier on \(Y\) with \(Y\neq X\) and \(Y\neq U\), the substitution
passes inside and \(\sem U_{s[Y\mapsto R]}=\sem U_s\), so the induction
hypothesis applies to each modified assignment.  Under a binder for \(X\),
\(X\) is not free and (L1) applies.  Under a binder for \(Y=U\), the
substitution first renames the binder to a fresh \(Y'\), which preserves
truth by the renaming fact, and then proceeds as before.  The renaming facts
are the instances \(U:=Y'\) of (L3), with \(\psi\) of smaller \(c\), together with
(L1), because \(Y'\) is fresh.  (L2) is proved in the same pattern with
object terms: atoms and equations by the term clause
\(\sem{r[t/x]}_s=\sem r_{s[x\mapsto\sem t_s]}\), proved by induction on \(r\);
relation binders are passed unchanged, and an object binder for a variable
of \(t\) is renamed first.

\emph{Soundness.}  We show that every rule carries
\(\mathsf H_{\mathsf{CC}}\)-valid premises to an \(\mathsf H_{\mathsf{CC}}\)-valid
conclusion and that every leaf is valid.  Id, \(\bot L\) and \({=}I\) are
valid, equality denoting identity.  Comp and Choice leaves are valid because
\(\mathsf H_{\mathsf{CC}}\)-models satisfy every eligible instance
(\cref{def:p02-sol-matched-henkin}); this is the only place where the
matched closure conditions are used.  Weakening, contraction, the
propositional rules, Cut and Mix are valid by the truth tables, applied to an
arbitrary model and assignment.  \(\forall L\) and \(\exists R\): by (L2),
\(\varphi[t/x]\) is true under \(s\) iff \(\varphi\) is true under
\(s[x\mapsto\sem t_s]\), which follows from \(\forall x\,\varphi\) and implies
\(\exists x\,\varphi\).  \(\forall R\): let the premise be valid, let \(s\)
satisfy \(\Gamma\) and falsify every member of \(\Delta\), and let \(a\in M\).
Since \(y\) is free in no formula of the conclusion, (L1) shows that
\(s[y\mapsto a]\) still satisfies \(\Gamma\) and falsifies \(\Delta\), so it
satisfies \(\varphi[y/x]\); by (L2), \(s[y\mapsto a][x\mapsto a]\) satisfies
\(\varphi\), and by (L1) so does \(s[x\mapsto a]\), because \(y\) is not free in
\(\forall x\,\varphi\).  Hence \(s\) satisfies \(\forall x\,\varphi\).
\(\exists L\) is the same argument on the left.  \({=}L\): if \(s\) satisfies
\(r=t\) and \(\varphi[t/x]\), then \(\sem r_s=\sem t_s\), and (L2) twice shows
that \(s\) satisfies \(\varphi[r/x]\).  \(\forall^2L\) and \(\exists^2R\): by
(L3), with \(\sem U_s\in\mathcal R_n\).  \(\forall^2R\) and \(\exists^2L\): the
argument for \(\forall R\) with \(R\in\mathcal R_n\) in place of \(a\), (L3) in
place of (L2), and the freshness of \(Y^n\).  By induction on derivations,
every natively derivable sequent is \(\mathsf H_{\mathsf{CC}}\)-valid.

\emph{Completeness: underivable pairs.}  Assume \(\Sigma\Rightarrow\Pi\) is not
natively derivable.  Call a pair \((\Delta,\Theta)\) of sets of formulas
\emph{underivable} if no sequent \(\Sigma'\Rightarrow\Pi'\) is derivable whose
antecedent members all lie in \(\Delta\) and whose succedent members all lie
in \(\Theta\).  By weakening and contraction, derivability of such a sequent
depends only on the two sets of members, so \((\Sigma,\Pi)\), read as sets,
is underivable.  Three facts hold for an underivable \((\Delta,\Theta)\).
\begin{enumerate}
 \item[(P1)] For every \(\varphi\), \((\Delta\cup\{\varphi\},\Theta)\) or
   \((\Delta,\Theta\cup\{\varphi\})\) is underivable.  Otherwise
   \(\Sigma_1,\varphi\Rightarrow\Pi_1\) and \(\Sigma_2\Rightarrow\Pi_2,\varphi\) are
   derivable with members in \(\Delta\) and \(\Theta\); Cut gives
   \(\Sigma_2,\Sigma_1\Rightarrow\Pi_2,\Pi_1\).
 \item[(P2)] If \(\exists x\,\varphi\in\Delta\) and \(y\) is free in no member of
   \(\Delta\cup\Theta\), then \((\Delta\cup\{\varphi[y/x]\},\Theta)\) is
   underivable: from a derivable \(\Sigma',\varphi[y/x]\Rightarrow\Pi'\),
   \(\exists L\) gives \(\Sigma',\exists x\,\varphi\Rightarrow\Pi'\), its freshness
   condition holding.  Likewise for \(\forall x\,\varphi\in\Theta\) with
   \(\forall R\), adding \(\varphi[y/x]\) to \(\Theta\); for
   \(\exists X^n\varphi\in\Delta\) with \(\exists^2L\); and for
   \(\forall X^n\varphi\in\Theta\) with \(\forall^2R\), using a relation
   variable \(Y^n\) free in no member of \(\Delta\cup\Theta\).
 \item[(P3)] If \(\Delta\cap\Theta\neq\varnothing\) or \(\bot\in\Delta\), the
   pair is derivable, by Id and \(\bot L\).
\end{enumerate}

\emph{Saturation.}  The formulas are finite strings over a countable
alphabet; enumerate them as \(\varphi_0,\varphi_1,\ldots\).  Put
\((\Delta_0,\Theta_0)=(\Sigma,\Pi)\).  Given the underivable pair of finite sets
\((\Delta_k,\Theta_k)\), use (P1) to add \(\varphi_k\) to \(\Delta_k\) if that
leaves the pair underivable and to \(\Theta_k\) otherwise.  If \(\varphi_k\)
went to the antecedent and is \(\exists x\,\psi\) or \(\exists X^n\psi\), or went
to the succedent and is \(\forall x\,\psi\) or \(\forall X^n\psi\), add the
instance of (P2) with the first object or \(n\)-ary relation variable free in
no member of the current finite sets; one exists, because each set is
finite and each variable stock is infinite.  Let \(\Delta^*\) and
\(\Theta^*\) be the unions.  The pair \((\Delta^*,\Theta^*)\) is underivable,
because a derivable sequent has finitely many members, all present at some
stage.  By construction and (P3), every formula lies in exactly one of
\(\Delta^*\), \(\Theta^*\).  Write \(\varphi\in\Delta^*\) as ``\(\varphi\) is in''.
For each of the following, the failure of the stated equivalence would put
into the pair a derivable sequent from the list after it, each derivable by
Id with one or two rules:
\begin{enumerate}
 \item[(H1)] \(\neg\varphi\) is in iff \(\varphi\) is not; \(\varphi\land\psi\) is in
   iff both are; \(\varphi\lor\psi\) is in iff one is; \(\varphi\to\psi\) is in iff
   \(\varphi\) is not or \(\psi\) is; \(\bot\) is not in.  Sequents:
   \(\neg\varphi,\varphi\Rightarrow\); \(\Rightarrow\varphi,\neg\varphi\);
   \(\varphi\land\psi\Rightarrow\varphi\); \(\varphi\land\psi\Rightarrow\psi\);
   \(\varphi,\psi\Rightarrow\varphi\land\psi\); \(\varphi\Rightarrow\varphi\lor\psi\);
   \(\psi\Rightarrow\varphi\lor\psi\); \(\varphi\lor\psi\Rightarrow\varphi,\psi\);
   \(\varphi\to\psi,\varphi\Rightarrow\psi\); \(\Rightarrow\varphi,\varphi\to\psi\);
   \(\psi\Rightarrow\varphi\to\psi\); \(\bot\Rightarrow\).
 \item[(H2)] \(\forall x\,\varphi\) is in iff \(\varphi[t/x]\) is in for every
   object term \(t\); \(\exists x\,\varphi\) is in iff \(\varphi[t/x]\) is in for
   some \(t\).  Sequents: \(\forall x\,\varphi\Rightarrow\varphi[t/x]\) and
   \(\varphi[t/x]\Rightarrow\exists x\,\varphi\); the other directions are the
   witnesses added by (P2).
 \item[(H3)] \(\forall X^n\varphi\) is in iff \(\varphi[U/X]\) is in for every
   \(n\)-ary relation variable or predicate constant \(U\); \(\exists X^n\varphi\)
   is in iff \(\varphi[U/X]\) is in for some such \(U\).  Sequents:
   \(\forall X^n\varphi\Rightarrow\varphi[U/X]\) and
   \(\varphi[U/X]\Rightarrow\exists X^n\varphi\) (\(\forall^2L\), \(\exists^2R\));
   the other directions are the (P2) witnesses.
 \item[(H4)] \(t=t\) is in; and if \(r=t\) and \(\varphi[r/x]\) are in, then
   \(\varphi[t/x]\) is in.  Sequents: \(\Rightarrow t=t\); and
   \(r=t,\varphi[r/x]\Rightarrow\varphi[t/x]\), derived as follows.  From
   \(\varphi[r/x]\Rightarrow\varphi[r/x]\), \(\neg L\) gives
   \(\neg\varphi[r/x],\varphi[r/x]\Rightarrow\); \({=}L\), applied to the formula
   \((\neg\varphi)[r/x]\), gives \(\varphi[r/x],r=t,\neg\varphi[t/x]\Rightarrow\);
   Cut with \(\Rightarrow\varphi[t/x],\neg\varphi[t/x]\) (Id and \(\neg R\)) gives
   the sequent.
 \item[(H5)] Every eligible Comp and Choice instance is in: \(\Rightarrow\sigma\)
   is a leaf.
\end{enumerate}

\emph{The canonical structure.}  For object terms put \(r\approx t\) iff
\(r=t\) is in.  By (H4), \(\approx\) is reflexive; taking
\(\varphi:=(x=r)\) with \(x\) not in \(r\) gives symmetry, and
\(\varphi:=(r'=x)\) gives transitivity.  It is a congruence: if \(r\approx t\),
then \(\varphi:=(f(\bar u,r,\bar v)=f(\bar u,x,\bar v))\) gives
\(f(\bar u,r,\bar v)\approx f(\bar u,t,\bar v)\), and
\(\varphi:=U(\bar u,x,\bar v)\) shows that \(U(\bar u,r,\bar v)\) is in iff
\(U(\bar u,t,\bar v)\) is in (the converse by symmetry).  Let \(M\) be the
set of \(\approx\)-classes \([t]\), interpret every object constant and
function symbol on classes of terms, and for every \(n\)-ary relation
variable or predicate constant \(U\) put
\[
 \operatorname{ext}(U)=\{([t_1],\ldots,[t_n]):U(t_1,\ldots,t_n)\text{ is in}\},
 \qquad
 \mathcal R_n=\{\operatorname{ext}(U):U\text{ an }n\text{-ary relation variable or predicate constant}\}.
\]
All are well defined by the congruence, \(M\) and each \(\mathcal R_n\) are
nonempty, and with \(P^{\mathcal H}=\operatorname{ext}(P)\in\mathcal R_n\) this is a
relational Henkin structure \(\mathcal H\).  Let \(s_0(x)=[x]\) and
\(s_0(X)=\operatorname{ext}(X)\), a matching assignment.  By induction on terms,
\(\sem t_{s_0}=[t]\), and \(\sem U_{s_0}=\operatorname{ext}(U)\).

\emph{Truth lemma.}  For every formula \(\varphi\),
\(\mathcal H,s_0\entails\varphi\) iff \(\varphi\) is in.  Induction on
\(c(\varphi)\).  Atoms: \(U(\bar t)\) holds iff \(([\bar t])\in\operatorname{ext}(U)\)
iff \(U(\bar t)\) is in; \(r=t\) holds iff \([r]=[t]\) iff \(r=t\) is in; \(\bot\)
by (H1).  Connectives: (H1).  \(\forall x\,\varphi\) holds iff
\(\mathcal H,s_0[x\mapsto[t]]\entails\varphi\) for every term \(t\), iff (L2)
\(\mathcal H,s_0\entails\varphi[t/x]\) for every \(t\), iff (induction) every
\(\varphi[t/x]\) is in, iff (H2) \(\forall x\,\varphi\) is in; \(\exists\) is the
same.  \(\forall X^n\varphi\) holds iff
\(\mathcal H,s_0[X\mapsto\operatorname{ext}(U)]\entails\varphi\) for every \(U\), iff
(L3) every \(\varphi[U/X]\) holds, iff (induction) every \(\varphi[U/X]\) is
in, iff (H3) \(\forall X^n\varphi\) is in; \(\exists X^n\) is the same.

\emph{Conclusion.}  By (H5) and the truth lemma, \(\mathcal H\) satisfies every
eligible Comp and Choice instance; these instances are sentences, since
\(\operatorname{ucl}\) closes all their parameters.  So \(\mathcal H\) is an
\(\mathsf H_{\mathsf{CC}}\)-model, \(s_0\) satisfies every member of
\(\Sigma\subseteq\Delta^*\) and none of \(\Pi\subseteq\Theta^*\), and
\(\Sigma\Rightarrow\Pi\) is not \(\mathsf H_{\mathsf{CC}}\)-valid.  The
enumeration, the stages, the quotient and \(s_0\) are explicitly defined, so
no form of the Axiom of Choice is used.
\end{proof}
\end{annexitem}

\begin{annexitem}{annex:p02-w02-p02-040-sol-correspondence}{thm:p02-w02-p02-040-sol-correspondence}{Native and translation-defined derivability coincide}
\begin{proof}
If \(\Gamma\derives_{\mathrm{nat}}\varphi\), some finite \(\Sigma\subseteq\Gamma\)
has \(\Sigma\Rightarrow\varphi\) natively derivable, so
\(\Sigma\entails_{\mathrm H}^{\mathsf{CC}}\varphi\) by
\cref{thm:p02-w02-p02-040-sol-native}, hence
\(\Gamma\entails_{\mathrm H}^{\mathsf{CC}}\varphi\), hence
\(\Gamma\derives_{\mathsf{SOL}_{\mathrm H}^{\mathsf{CC}}}\varphi\) by
\cref{thm:p02-sol-derived-adequacy}.  Conversely, if
\(\Gamma\derives_{\mathsf{SOL}_{\mathrm H}^{\mathsf{CC}}}\varphi\), then by
\cref{def:p02-sol-derived-calculus} a natural-deduction derivation of
\(\varphi^\sharp\) exists whose open assumptions, finitely many, lie in the
translated background set; so
\(\Sigma\derives_{\mathsf{SOL}_{\mathrm H}^{\mathsf{CC}}}\varphi\) for a finite
\(\Sigma\subseteq\Gamma\).  By \cref{thm:p02-sol-derived-adequacy},
\(\Sigma\entails_{\mathrm H}^{\mathsf{CC}}\varphi\), and by
\cref{thm:p02-w02-p02-040-sol-native}, \(\Sigma\Rightarrow\varphi\) is natively
derivable, so \(\Gamma\derives_{\mathrm{nat}}\varphi\).  The third relation equals
the second by \cref{thm:p02-sol-derived-adequacy}.

For a finite sequent \(\Sigma\Rightarrow\Pi\), translation-defined
derivability is
\(T^\sharp_{\mathsf{CC}}\cup\Sigma^\sharp\derives_{\mathsf{MSFOL}}\bigvee\Pi^\sharp\)
(\cref{def:p02-sol-derived-calculus}).  The translation commutes with the
connectives and \(\bot^\sharp=\bot\), so \(\bigvee\Pi^\sharp\) is
\((\bigvee\Pi)^\sharp\) for the same listing, and the condition is
\(\Sigma\derives_{\mathsf{SOL}_{\mathrm H}^{\mathsf{CC}}}\bigvee\Pi\).  By
\cref{thm:p02-sol-derived-adequacy} this holds iff
\(\Sigma\entails_{\mathrm H}^{\mathsf{CC}}\bigvee\Pi\), that is, iff the sequent is
\(\mathsf H_{\mathsf{CC}}\)-valid, and by \cref{thm:p02-w02-p02-040-sol-native}
iff it is natively derivable.
\end{proof}
\end{annexitem}

\begin{annexitem}{annex:p02-w02-p02-040-pa2-categorical}{thm:p02-w02-p02-040-pa2-categorical}{Full categoricity of \(\mathsf{PA}_2\) with addition and multiplication}
\begin{proof}
\emph{The standard structure is a full model.}  In \(\mathcal N\), \(n+1\neq0\)
and \(n+1=m+1\) implies \(n=m\).  For the induction conjunct, let
\(R\subseteq\N\) be arbitrary (full semantics) with \(0\in R\) and
\(n+1\in R\) whenever \(n\in R\); induction in the metatheory gives
\(R=\N\).  The four recursion conjuncts are the recursion equations that
define \(+\) and \(\times\) on \(\N\).

\emph{The map is a bijection preserving \(0\) and \(S\).}  The conjuncts
\(\forall x\,(Sx\neq0)\), \(\forall x\,\forall y\,(Sx=Sy\to x=y)\) and the
induction sentence are three of the seven conjuncts of \(\mathsf{PA}_2\),
so the \(\{0,S\}\)-reduct of \(\mathcal M\) satisfies the displayed
second-order Peano axioms of \cref{sec:second-order-logic-03} under full
semantics.  The proof of the categoricity theorem there shows that \(f\) is
injective and surjective, and \(f\) preserves \(0\) and \(S\) by its
definition.

\emph{Addition.}  Fix \(m\) and argue by induction on \(n\).  For \(n=0\),
\(f(m+0)=f(m)=f(m)+^{\mathcal M}0^{\mathcal M}=f(m)+^{\mathcal M}f(0)\) by
the conjunct \(\forall x\,(x+0=x)\).  If \(f(m+n)=f(m)+^{\mathcal M}f(n)\),
then
\[
 f(m+(n+1))=S^{\mathcal M}(f(m+n))
 =S^{\mathcal M}(f(m)+^{\mathcal M}f(n))
 =f(m)+^{\mathcal M}S^{\mathcal M}(f(n))
 =f(m)+^{\mathcal M}f(n+1),
\]
using the definition of \(f\), the induction hypothesis and the conjunct
\(\forall x\,\forall y\,(x+S(y)=S(x+y))\).

\emph{Multiplication.}  Fix \(m\).  For \(n=0\),
\(f(m\times0)=f(0)=0^{\mathcal M}=f(m)\times^{\mathcal M}0^{\mathcal M}\) by
\(\forall x\,(x\times0=0)\).  If \(f(m\times n)=f(m)\times^{\mathcal M}f(n)\),
then, using the addition clause just proved,
\[
 f(m\times(n+1))=f((m\times n)+m)
 =(f(m)\times^{\mathcal M}f(n))+^{\mathcal M}f(m)
 =f(m)\times^{\mathcal M}S^{\mathcal M}(f(n))
 =f(m)\times^{\mathcal M}f(n+1)
\]
by \(\forall x\,\forall y\,(x\times S(y)=(x\times y)+x)\).

So \(f\) is an isomorphism of \(\{0,S,+,\times\}\)-structures.  If
\(\mathcal M_1,\mathcal M_2\) are full models with maps \(f_1,f_2\), then
\(f_2\circ f_1^{-1}\) is an isomorphism from \(\mathcal M_1\) onto
\(\mathcal M_2\).
\end{proof}
\end{annexitem}

\begin{annexitem}{annex:p02-w02-p02-040-reals-categorical}{thm:p02-w02-p02-040-reals-categorical}{Full categoricity and uncountability of complete ordered fields}
\begin{proof}
Let \(\mathcal F\) be a complete ordered field with domain \(F\).  Under
full semantics the completeness sentence applies to every subset of \(F\):
every nonempty subset with an upper bound has a least upper bound, written
\(\sup_{\mathcal F}\).

\emph{Step 1: arithmetic consequences of the axioms.}  Additive inverses
are unique: if \(x+y=0=x+y'\), then \(y=y+(x+y')=(y+x)+y'=y'\) by
F1--F3; write \(-x\) and \(x-y\defeq x+(-y)\).  A nonzero \(x\) has an inverse by F8, and it is
unique: if \(xy=1=xy'\), then \(y=y(xy')=(yx)y'=y'\) by F5--F7; write
\(x^{-1}\).  From \(x0=x(0+0)=x0+x0\) (F3, F9),
adding \(-(x0)\) gives \(x0=0\).  From
\(x(-1)+x=x((-1)+1)=x0=0\) we get \((-1)x=-x\), and then
\((-1)(-1)=-(-1)=1\).  For the order: by O4, \(x<y\) iff \(0<y-x\).  If
\(x<y\) and \(0<z\), then \(0<(y-x)z=yz-xz\) by O5 and F9, so \(xz<yz\).
If \(1<0\), adding \(-1\) gives \(0<-1\), so \(0<(-1)(-1)=1\) by O5,
and O2 with O1 gives a contradiction; since \(0\neq1\) (F10), O3 gives
\(0<1\).  If \(0<x\) then \(0<x^{-1}\): otherwise \(x^{-1}<0\) (it is not
\(0\), since \(x x^{-1}=1\neq0=x0\)), so \(0<-x^{-1}\) and
\(0<x(-x^{-1})=-1\), contradicting \(0<1\).  The relation \(\le\) is a
total order by O1--O3.

\emph{Step 2: the rationals inside \(\mathcal F\).}  First some further
identities.  Since \(x+(-x)=0\), \(x\) is the additive inverse of \(-x\), so
\(-(-x)=x\) by uniqueness.  Since \((x+y)+((-x)+(-y))=(x+(-x))+(y+(-y))=0\) by
F1 and F2, \(-(x+y)=(-x)+(-y)\).  By F5, F6 and Step~1,
\((-x)y=((-1)x)y=(-1)(xy)=-(xy)\), likewise \(x(-y)=-(xy)\), and
\((-x)(-y)=-(x(-y))=-(-(xy))=xy\).  If \(z\neq0\) and \(xz=yz\), then
multiplying by \(z^{-1}\) and using F5, F7 gives \(x=y\).  If \(0<y\), adding
\(-y\) to both sides (O4) gives \(-y<0\).

Define \(n_{\mathcal F}\) for \(n\in\N\) by \(0_{\mathcal F}=0\) and
\((n+1)_{\mathcal F}=n_{\mathcal F}+1\).  Induction on \(n\), adding
\(n_{\mathcal F}\) to \(0<1\) by O4 and using O2, gives
\(0<n_{\mathcal F}<(n+1)_{\mathcal F}\) for \(n\geq1\); in particular
\(n_{\mathcal F}\neq0\) for \(n\ge1\).
\begin{enumerate}
 \item[(a)] \((m+n)_{\mathcal F}=m_{\mathcal F}+n_{\mathcal F}\) for
   \(m,n\in\N\), by induction on \(n\): \((m+0)_{\mathcal F}=m_{\mathcal F}=m_{\mathcal F}+0\)
   by F3, and \((m+n+1)_{\mathcal F}=(m+n)_{\mathcal F}+1=(m_{\mathcal F}+n_{\mathcal F})+1
   =m_{\mathcal F}+(n_{\mathcal F}+1)=m_{\mathcal F}+(n+1)_{\mathcal F}\) by F1.
 \item[(b)] \((mn)_{\mathcal F}=m_{\mathcal F}n_{\mathcal F}\) for \(m,n\in\N\),
   by induction on \(n\): \((m0)_{\mathcal F}=0=m_{\mathcal F}0\) by Step~1,
   and \((m(n+1))_{\mathcal F}=(mn+m)_{\mathcal F}=m_{\mathcal F}n_{\mathcal F}+m_{\mathcal F}1
   =m_{\mathcal F}(n_{\mathcal F}+1)=m_{\mathcal F}(n+1)_{\mathcal F}\) by (a), F7 and F9.
 \item[(c)] \((m-n)_{\mathcal F}=m_{\mathcal F}-n_{\mathcal F}\) for
   \(m\ge n\) in \(\N\): by (a), \((m-n)_{\mathcal F}+n_{\mathcal F}=m_{\mathcal F}\);
   add \(-n_{\mathcal F}\).
\end{enumerate}
For a negative integer \(k\) put \(k_{\mathcal F}=-((-k)_{\mathcal F})\).
Then \((-k)_{\mathcal F}=-(k_{\mathcal F})\) for every integer \(k\): for
\(k>0\) by definition, for \(k<0\) because
\(-(k_{\mathcal F})=-(-((-k)_{\mathcal F}))=(-k)_{\mathcal F}\), and for \(k=0\)
because \(-0=0\).
\begin{enumerate}
 \item[(d)] \((k+l)_{\mathcal F}=k_{\mathcal F}+l_{\mathcal F}\) for all integers.
   If \(k,l\ge0\), this is (a).  If \(k,l\le0\),
   \((k+l)_{\mathcal F}=-((-k)+(-l))_{\mathcal F}=-((-k)_{\mathcal F}+(-l)_{\mathcal F})
   =-((-k)_{\mathcal F})+(-((-l)_{\mathcal F}))=k_{\mathcal F}+l_{\mathcal F}\).
   If \(k\ge0>l\) and \(k+l\ge0\), then by (c)
   \((k+l)_{\mathcal F}=(k-(-l))_{\mathcal F}=k_{\mathcal F}-(-l)_{\mathcal F}=k_{\mathcal F}+l_{\mathcal F}\).
   If \(k\ge0>l\) and \(k+l<0\), then \(-l>k\), and by (c)
   \((k+l)_{\mathcal F}=-((-l)-k)_{\mathcal F}=-((-l)_{\mathcal F}-k_{\mathcal F})
   =-((-l)_{\mathcal F})+k_{\mathcal F}=l_{\mathcal F}+k_{\mathcal F}\), using
   \(-(x-y)=(-x)+y\) and F2.  The case \(k<0\le l\) follows by F2.
 \item[(e)] \((kl)_{\mathcal F}=k_{\mathcal F}l_{\mathcal F}\) for all integers.
   If \(k,l\ge0\), this is (b).  If \(k=-a\) with \(a>0\) and \(l\ge0\), then
   \((kl)_{\mathcal F}=-((al)_{\mathcal F})=-(a_{\mathcal F}l_{\mathcal F})=(-a_{\mathcal F})l_{\mathcal F}
   =k_{\mathcal F}l_{\mathcal F}\); the case \(k\ge0>l\) is symmetric by F6.
   If \(k=-a\), \(l=-b\) with \(a,b>0\), then
   \((kl)_{\mathcal F}=(ab)_{\mathcal F}=a_{\mathcal F}b_{\mathcal F}=(-a_{\mathcal F})(-b_{\mathcal F})=k_{\mathcal F}l_{\mathcal F}\).
 \item[(f)] \(k<l\) iff \(k_{\mathcal F}<l_{\mathcal F}\).  By (d),
   \(l_{\mathcal F}-k_{\mathcal F}=(l-k)_{\mathcal F}\), and
   \(j_{\mathcal F}>0\) iff \(j>0\): positive \(j\) give \(0<j_{\mathcal F}\),
   \(0_{\mathcal F}=0\), and negative \(j\) give \(j_{\mathcal F}=-((-j)_{\mathcal F})<0\).
   By Step~1, \(k_{\mathcal F}<l_{\mathcal F}\) iff \(0<l_{\mathcal F}-k_{\mathcal F}\).
\end{enumerate}
For \(q=k/n\) with \(k\in\mathbb Z\) and
\(n\geq1\) put \(q_{\mathcal F}=k_{\mathcal F}(n_{\mathcal F})^{-1}\).  This
does not depend on the representation: \(k/n=k'/n'\) means \(kn'=k'n\), so
\(k_{\mathcal F}n'_{\mathcal F}=k'_{\mathcal F}n_{\mathcal F}\) by (e), and
multiplying by \((n_{\mathcal F}n'_{\mathcal F})^{-1}\) gives equality.  For
\(q=k/n\) and \(q'=k'/n'\), put \(z=n_{\mathcal F}n'_{\mathcal F}\), which is
positive by O5 and so nonzero.  Since \(q+q'=(kn'+k'n)/(nn')\), (d) and (e) give
\((q+q')_{\mathcal F}z=k_{\mathcal F}n'_{\mathcal F}+k'_{\mathcal F}n_{\mathcal F}\),
the inverse of \(z\) cancelling by F5 and F7; and by F9, F5 and F6,
\((q_{\mathcal F}+q'_{\mathcal F})z=k_{\mathcal F}n_{\mathcal F}^{-1}n_{\mathcal F}n'_{\mathcal F}
+k'_{\mathcal F}n'^{-1}_{\mathcal F}n_{\mathcal F}n'_{\mathcal F}=k_{\mathcal F}n'_{\mathcal F}+k'_{\mathcal F}n_{\mathcal F}\).
Cancelling \(z\) gives \((q+q')_{\mathcal F}=q_{\mathcal F}+q'_{\mathcal F}\).
Likewise \(qq'=(kk')/(nn')\) gives
\((qq')_{\mathcal F}z=k_{\mathcal F}k'_{\mathcal F}=(q_{\mathcal F}q'_{\mathcal F})z\),
so \((qq')_{\mathcal F}=q_{\mathcal F}q'_{\mathcal F}\).  Finally
\(q<q'\) iff \(q_{\mathcal F}<q'_{\mathcal F}\): with common denominator
\(n\geq1\), \(k/n<k'/n\) iff \(k<k'\) iff \(k_{\mathcal F}<k'_{\mathcal F}\) iff
\(k_{\mathcal F}n_{\mathcal F}^{-1}<k'_{\mathcal F}n_{\mathcal F}^{-1}\), the
last step multiplying by the positive \(n_{\mathcal F}^{-1}\) or by
\(n_{\mathcal F}\).

\emph{Step 3: Archimedes.}  For every \(x\in F\) some \(n\in\N\) has
\(x<n_{\mathcal F}\).  Otherwise \(\{n_{\mathcal F}:n\in\N\}\) is nonempty and
bounded above by \(x\); let \(s\) be its supremum.  Since \(s-1<s\), the
element \(s-1\) is not an upper bound, so \(s-1<n_{\mathcal F}\) for some
\(n\), whence \(s<(n+1)_{\mathcal F}\), contradicting that \(s\) is an upper
bound.  Applying this to \(-x\) gives an integer \(k\) with
\(k_{\mathcal F}<x\).

\emph{Step 4: density.}  If \(a<b\) in \(F\), some \(q\in\mathbb Q\) has
\(a<q_{\mathcal F}<b\).  By Step~3 choose \(n\geq1\) with
\((b-a)^{-1}<n_{\mathcal F}\), so \(1<n_{\mathcal F}(b-a)\).  The set of
integers \(k\) with \(n_{\mathcal F}a<k_{\mathcal F}\) is nonempty and bounded
below, by Step~3; let \(k\) be its least member.  Then
\((k-1)_{\mathcal F}\le n_{\mathcal F}a\), so
\(k_{\mathcal F}\le n_{\mathcal F}a+1<n_{\mathcal F}a+n_{\mathcal F}(b-a)=n_{\mathcal F}b\).
Multiplying \(n_{\mathcal F}a<k_{\mathcal F}<n_{\mathcal F}b\) by
\(n_{\mathcal F}^{-1}>0\) gives \(a<(k/n)_{\mathcal F}<b\).

\emph{Step 5: cuts determine elements.}  For \(x\in F\) the cut
\(L_{\mathcal F}(x)\) is nonempty and bounded above in \(\mathbb Q\) (Step~3),
and \(x=\sup_{\mathcal F}\{q_{\mathcal F}:q\in L_{\mathcal F}(x)\}\): \(x\) is
an upper bound, and if \(y<x\) were one, Step~4 would give \(q\) with
\(y<q_{\mathcal F}<x\).  Hence \(x<y\) iff \(L_{\mathcal F}(x)\subsetneq L_{\mathcal F}(y)\):
if \(x<y\), clearly \(L_{\mathcal F}(x)\subseteq L_{\mathcal F}(y)\), and Step~4
gives \(q\) with \(x<q_{\mathcal F}<y\), which lies in the second set only;
the converse follows from O3 and what was just shown.  In particular
\(L_{\mathcal F}(x)=L_{\mathcal F}(y)\) implies \(x=y\).

\emph{Step 6: the map.}  Let \(\mathcal G\) be a second complete ordered
field.  For \(x\in F\), the set \(\{q_{\mathcal G}:q\in L_{\mathcal F}(x)\}\)
is nonempty and bounded above in \(G\): if \(x<m_{\mathcal F}\) then every
\(q\in L_{\mathcal F}(x)\) has \(q<m\), so \(q_{\mathcal G}<m_{\mathcal G}\).  So
\(h(x)\) is defined.  We show \(L_{\mathcal G}(h(x))=L_{\mathcal F}(x)\).  If
\(q\in L_{\mathcal F}(x)\), Step~4 gives \(r\) with \(q_{\mathcal F}<r_{\mathcal F}<x\),
so \(q_{\mathcal G}<r_{\mathcal G}\le h(x)\).  If \(q_{\mathcal G}<h(x)\), then
\(q_{\mathcal G}\) is not an upper bound, so \(q_{\mathcal G}<r_{\mathcal G}\) for
some \(r\in L_{\mathcal F}(x)\); then \(q<r\) and \(q_{\mathcal F}<r_{\mathcal F}<x\).
By Step~5 in \(\mathcal G\), \(h\) is the only map with this property.

\emph{Step 7: order, bijectivity, rationals.}  By Steps~5 and~6,
\(x<y\) iff \(L_{\mathcal F}(x)\subsetneq L_{\mathcal F}(y)\) iff
\(L_{\mathcal G}(h(x))\subsetneq L_{\mathcal G}(h(y))\) iff \(h(x)<h(y)\); so
\(h\) preserves and reflects \(<\) and is injective.  The same construction
from \(\mathcal G\) to \(\mathcal F\) gives \(k\) with
\(L_{\mathcal F}(k(y))=L_{\mathcal G}(y)\); then
\(L_{\mathcal G}(h(k(y)))=L_{\mathcal G}(y)\), so \(h(k(y))=y\) and \(h\) is
onto.  Since \(L_{\mathcal F}(q_{\mathcal F})=\{r:r<q\}=L_{\mathcal G}(q_{\mathcal G})\),
\(h(q_{\mathcal F})=q_{\mathcal G}\); in particular \(h(0)=0\) and \(h(1)=1\).

\emph{Step 8: addition.}  For \(x,y\in F\),
\[
 L_{\mathcal F}(x+y)=\{r\in\mathbb Q:\ r<p+p'\text{ for some }p\in L_{\mathcal F}(x),\ p'\in L_{\mathcal F}(y)\}.
\]
If \(r_{\mathcal F}<x+y\), then \(r_{\mathcal F}-y<x\); Step~4 gives \(p\) with
\(r_{\mathcal F}-y<p_{\mathcal F}<x\), hence \(r_{\mathcal F}-p_{\mathcal F}<y\), and
gives \(p'\) with \(r_{\mathcal F}-p_{\mathcal F}<p'_{\mathcal F}<y\); then
\(r<p+p'\) by Step~2.  Conversely \(r<p+p'\) with \(p_{\mathcal F}<x\),
\(p'_{\mathcal F}<y\) gives \(r_{\mathcal F}<p_{\mathcal F}+p'_{\mathcal F}<x+y\) by
O4 twice.  The same identity holds in \(\mathcal G\).  Applying it to
\(h(x),h(y)\) and using Step~6,
\(L_{\mathcal G}(h(x)+h(y))\) is the displayed set, which is
\(L_{\mathcal F}(x+y)=L_{\mathcal G}(h(x+y))\); Step~5 gives
\(h(x+y)=h(x)+h(y)\).  Hence also \(h(-x)=-h(x)\).

\emph{Step 9: multiplication.}  For \(0<x\) and \(0<y\),
\[
 L_{\mathcal F}(xy)=\{r:r\le0\}\cup
 \{r:\ r<pp'\text{ for some }p\in L_{\mathcal F}(x),\ p'\in L_{\mathcal F}(y)\text{ with }0<p,\ 0<p'\}.
\]
Indeed, if \(0<r_{\mathcal F}<xy\), then \(0<r_{\mathcal F}y^{-1}<x\); Step~4
gives \(p\) with \(r_{\mathcal F}y^{-1}<p_{\mathcal F}<x\), so \(0<p\) and
\(r_{\mathcal F}p_{\mathcal F}^{-1}<y\), and gives \(p'\) with
\(r_{\mathcal F}p_{\mathcal F}^{-1}<p'_{\mathcal F}<y\), so \(0<p'\) and \(r<pp'\).
Conversely, \(r\le0\) gives \(r_{\mathcal F}\le0<xy\) by O5, and
\(0<p_{\mathcal F}<x\), \(0<p'_{\mathcal F}<y\), \(r<pp'\) give
\(r_{\mathcal F}<p_{\mathcal F}p'_{\mathcal F}<xp'_{\mathcal F}<xy\) by Step~1.  The
same identity holds in \(\mathcal G\); since \(h\) preserves positivity
(Step~7), the argument of Step~8 gives \(h(xy)=h(x)h(y)\) for positive
\(x,y\).  The remaining cases follow from \(x0=0\), \(h(0)=0\),
\(h(-x)=-h(x)\) and \((-x)y=-(xy)\), since every nonzero element is
\(z\) or \(-z\) with \(0<z\).

So \(h\) is an isomorphism of \(\Lang_{\mathsf R}\)-structures from
\(\mathcal F\) onto \(\mathcal G\).

\emph{Step 10: uncountability.}  Let \((r_n)_{n\in\N}\) be any sequence in
\(F\).  Define \(a_0=0\), \(b_0=1\) and, given \(a_n<b_n\), put
\(c=a_n+(b_n-a_n)3_{\mathcal F}^{-1}\) and \(d=a_n+(b_n-a_n)2_{\mathcal F}3_{\mathcal F}^{-1}\),
so \(a_n<c<d<b_n\) by Step~1; let \((a_{n+1},b_{n+1})=(d,b_n)\) if
\(r_n\le c\) and \((a_n,c)\) otherwise.  This recursion makes no choices.
In either case \(r_n\notin[a_{n+1},b_{n+1}]\),
\(a_n\le a_{n+1}<b_{n+1}\le b_n\).  By induction \(a_m<b_n\) for all \(m,n\):
for \(m\le n\), \(a_m\le a_n<b_n\); for \(m>n\), \(a_m<b_m\le b_n\).  Let
\(s=\sup_{\mathcal F}\{a_m:m\in\N\}\), which exists because \(b_0\) is an
upper bound.  Then \(a_{n+1}\le s\le b_{n+1}\) for every \(n\), since each
\(b_{n+1}\) is an upper bound, so \(s\neq r_n\).  Thus no sequence lists all
of \(F\), and \(F\), being nonempty, is uncountable.
\end{proof}
\end{annexitem}

\begin{annexitem}{annex:p02-w02-p02-040-no-effective-calculus}{thm:p02-w02-p02-040-no-effective-calculus}{No effective complete calculus for full second-order validity}
\begin{proof}
\emph{Clause 1.}  The sentence \(\mathsf{PA}_2\) is fixed, so
\(p=\ulcorner\mathsf{PA}_2\urcorner\) is a fixed number, and
\(r(c)=C(10,C(p,C(c,0)))\) is a composition of the primitive-recursive
pairing \(C\) with constants.  By the constructor table of
\cref{sec:arithmetic-coding-04}, \(N_{10}(\ulcorner\varphi\urcorner,\ulcorner\psi\urcorner)\)
is the code of \(\varphi\to\psi\), and this clause is unchanged for
second-order formulas.  The section proves by a course-of-values
recursion that recognizing codes of first-order formulas and computing
their free-variable lists is primitive recursive.
\Cref{thm:p04-primitive-recursive-total} makes these numerical
functions computable, so the set of codes of first-order
\(\Lang_A\)-sentences is decidable.

\emph{Second-order recognition.}  Decidability of the set of codes of
second-order sentences is used in clauses 4 and 5; here are the clauses.
Keep the term test \(\mathrm{Tm}(c)\) and term variable lists \(\mathrm{V}(c)\)
of \cref{sec:arithmetic-coding-04}.  By one simultaneous course-of-values
recursion on \(c\) define the test \(\mathrm{Fm}_2(c)\), the list
\(\mathrm{FV}_1(c)\) of free individual-variable indices and the list
\(\mathrm{FV}_2(c)\) of free relation variables, each relation variable being
represented by its code \(N_{15}(n,i)\).  Write \(\mathrm{RV}(v)\) for: the
outer tag of \(v\) is \(15\), its argument list has length \(2\), and its first
entry \(n\) is at least \(1\).  At row \(c\), with \(k\) the outer tag and
\(a=t(c)\) the argument list:
\begin{enumerate}
 \item \(5\le k\le12\): the first-order formula clauses of the constructor table, with
   every formula child tested by \(\mathrm{Fm}_2\) at its earlier row instead
   of the first-order formula test; \(\mathrm{FV}_1\) as there (at
   \(N_{11},N_{12}\) remove the binder index); \(\mathrm{FV}_2(c)\) is empty
   for the atoms \(N_5\) and \(N_6\), and otherwise the union of
   \(\mathrm{FV}_2\) of the formula children.
 \item \(k=16\): \(\ell(a)=2\), \(v=(a)_0\) satisfies \(\mathrm{RV}(v)\), \(L=(a)_1\)
   has \(\ell(L)=(t(v))_0\) (arity matching), and \(\mathrm{Tm}((L)_j)\) for
   every \(j<\ell(L)\).  Then \(\mathrm{FV}_1(c)\) is the union of the
   \(\mathrm V((L)_j)\) and \(\mathrm{FV}_2(c)=[v]\).
 \item \(k=17\) or \(k=18\): \(\ell(a)=3\), \((a)_0\ge1\) and
   \(\mathrm{Fm}_2((a)_2)\).  Then \(\mathrm{FV}_1(c)=\mathrm{FV}_1((a)_2)\) and
   \(\mathrm{FV}_2(c)\) is \(\mathrm{FV}_2((a)_2)\) with every entry equal to
   \(N_{15}((a)_0,(a)_1)\) removed.
 \item Every other row, including \(c=0\), the term tags \(0\)--\(4\) and the tag \(15\), fails
   \(\mathrm{Fm}_2\), and its lists are empty.
\end{enumerate}
Every child consulted is a proper component of \(c\): the arguments, list
entries and list codes are all less than \(c\) by the property of \(C\)
recorded in \cref{sec:arithmetic-coding-04}.  Each clause is a bounded
combination of tag and length tests, list access, a loop of length
\(\ell(L)\), union and filtering, all primitive recursive there, so the
course-of-values principle of that section makes \(\mathrm{Fm}_2\),
\(\mathrm{FV}_1\) and \(\mathrm{FV}_2\) primitive recursive.  Induction on
\(c\) shows that \(\mathrm{Fm}_2(c)\) holds iff \(c\) is the code of a
second-order formula of \cref{def:p02-w02-p02-040-so-coding}, and that the
two lists are then its free individual variables and its free relation
variables: the four clauses are exactly the formation rules of
\cref{sec:second-order-logic-01} (an atom \(X^n_i(t_1,\ldots,t_n)\) needs
exactly \(n\) term arguments, and a second-order quantifier binds exactly the
relation variable of its arity and index), and codes are unique by
constructor injectivity.  The test
\(\mathrm{Sent}_2(c)\iff\mathrm{Fm}_2(c)\wedge\ell(\mathrm{FV}_1(c))=0\wedge\ell(\mathrm{FV}_2(c))=0\)
is primitive recursive, and by \cref{sec:second-order-logic-01} it holds iff
\(c\) codes a second-order sentence.  A code passing \(\mathrm{Fm}_2\) in which no
tag \(15\)--\(18\) occurs is a first-order formula code, so the first-order
test is unchanged.  The outer tags \(15\)--\(18\) are also used in
\cref{sec:arithmetic-coding-04} for derivation records, proof codes and
alpha-comparison records.  Those objects belong to other sorts and are never
submitted to \(\mathrm{Fm}_2\), and in any case they fail its clauses:
\(N_{15}(R)\), \(N_{17}(i)\) and \(N_{18}(j)\) have one argument rather than two or
three, and a proof code \(N_{16}(d,W)\) has \(d=N_{15}(R)\), whose argument list
has length one, so \(\mathrm{RV}(d)\) fails.

\emph{Isomorphism invariance.}  If \(f\) is an isomorphism from
\(\mathcal A\) onto \(\mathcal B\), then for every first-order formula
\(\varphi\) and assignment \(s\), \(\mathcal A,s\entails\varphi\) iff
\(\mathcal B,f\circ s\entails\varphi\).  Induction on terms gives
\(f(\sem{t}^{\mathcal A}_s)=\sem{t}^{\mathcal B}_{f\circ s}\); atoms use this
and injectivity of \(f\); connectives are immediate; and quantifiers use
that \(f\) is onto, since \(f\circ(s[x\mapsto a])=(f\circ s)[x\mapsto f(a)]\).

\emph{Clause 2.}  Suppose \(\entails_{\mathrm{full}}(\mathsf{PA}_2\to\sigma)\).
By \cref{thm:p02-w02-p02-040-pa2-categorical}, \(\mathcal N\) satisfies
\(\mathsf{PA}_2\) under full semantics, so \(\mathcal N\entails\sigma\).
Conversely suppose \(\mathcal N\entails\sigma\), and let \(\mathcal M\) be any
structure for \(\{0,S,+,\times\}\) and \(s\) any full assignment.  If
\(\mathcal M,s\) does not satisfy the sentence \(\mathsf{PA}_2\), the
implication holds.  If it does, \cref{thm:p02-w02-p02-040-pa2-categorical}
gives an isomorphism from \(\mathcal N\) onto \(\mathcal M\), and isomorphism
invariance gives \(\mathcal M\entails\sigma\).  So the implication holds in
every full structure.

\emph{Clause 3.} Suppose \(\mathrm{TA}\) were computably enumerable.
By the explicit halting-formula construction of
\cref{thm:p02-w02b-p02-034-ce-sigma-one}, it belongs to
\(\boldsymbol\Sigma^0_1\): some
\(\Sigma^0_1\) formula \(\mathsf{Tr}(x)\) satisfies
\(c\in\mathrm{TA}\) iff \(\mathcal N\entails\mathsf{Tr}(\bar c)\), for every
\(c\in\N\).  Then for every \(\Lang_A\)-sentence \(\sigma\),
\(\mathcal N\entails\mathsf{Tr}(\overline{\ulcorner\sigma\urcorner})\) iff
\(\ulcorner\sigma\urcorner\in\mathrm{TA}\) iff \(\mathcal N\entails\sigma\).  So
\(\mathcal N\entails\mathsf{Tr}(\overline{\ulcorner\sigma\urcorner})\leftrightarrow\sigma\)
for all \(\sigma\), and \(\mathsf{Tr}\) is an arithmetical truth definition,
contrary to \cref{thm:tarski-undefinability}.  This is the derivation of the
corollary stated after that theorem.

\emph{Clause 4.} Suppose \(\mathrm{Val}_2\) were the domain of a
partial computable \(\psi\). By
\cref{thm:p04-tm-partial-recursive}, \(\psi\) is partial recursive.
Let \(b(c)\) be the primitive-recursive test with value zero on
first-order \(\Lang_A\)-sentence codes and one on other numbers.
The function
\[
 g(c)=\mu z\,b(c)
\]
is defined with value zero exactly on those sentence codes. Its
minimized function is total and ignores \(z\), so the strict
minimization convention has precisely this effect. By strict
composition and addition,
\[
 h(c)=\psi(r(c))+g(c)
\]
is partial recursive and is defined exactly when \(c\) is such a
sentence code and \(r(c)\in\mathrm{Val}_2\). The same numerical
equivalence theorem supplies a machine computing \(h\), including
divergence outside this domain. Clause~2 identifies its domain with
\(\mathrm{TA}\), contradicting clause~3. Both translations preserve the fixed tuple encoding and
origin-based numerical output convention.

\emph{Clause 5.} If the theorem codes of such a \(\mathsf D\)
were the domain of a partial computable \(\chi\), repeat the
construction with the second-order sentence test and with \(r\)
the identity. It produces a partial computable function defined
exactly on the second-order sentences provable in \(\mathsf D\).
The assumed soundness-and-completeness biconditional makes this
set \(\mathrm{Val}_2\), contradicting clause~4.
\end{proof}
\end{annexitem}
